\documentclass[11pt]{article}
\usepackage[margin=1.05in]{geometry}
\usepackage{amsmath,amssymb,amsthm,mathtools,mathrsfs}
\usepackage{microtype}
\microtypesetup{expansion=false}
\usepackage{xcolor}
\definecolor{LinkBlue}{HTML}{1F4E79}
\definecolor{CiteGreen}{HTML}{2F6B3C}
\definecolor{UrlBrown}{HTML}{843C39}
\usepackage[colorlinks=true,linkcolor=LinkBlue,citecolor=CiteGreen,urlcolor=UrlBrown,breaklinks=true]{hyperref}
\newtheorem{theorem}{Theorem}[section]
\newtheorem{lemma}[theorem]{Lemma}
\newtheorem{proposition}[theorem]{Proposition}
\newtheorem{corollary}[theorem]{Corollary}
\theoremstyle{definition}
\newtheorem{definition}[theorem]{Definition}
\newtheorem{hypothesis}[theorem]{Hypothesis}
\theoremstyle{remark}
\newtheorem{remark}[theorem]{Remark}
\usepackage[nameinlink,capitalise,noabbrev]{cleveref}
\crefname{lemma}{Lemma}{Lemmas}
\crefname{proposition}{Proposition}{Propositions}
\crefname{corollary}{Corollary}{Corollaries}
\crefname{hypothesis}{Hypothesis}{Hypotheses}
\numberwithin{equation}{section}
\newcommand{\CH}{\operatorname{CH}}
\newcommand{\an}{\mathrm{an}}
\newcommand{\B}{\mathcal B}
\newcommand{\Ob}{\mathfrak o}
\newcommand{\Tr}{\operatorname{Tr}}
\newcommand{\Spec}{\operatorname{Spec}}
\newcommand{\id}{\operatorname{id}}
\newcommand{\pr}{\operatorname{pr}}
\newcommand{\BC}{\mathrm{BC,act}}
\newcommand{\homol}{\mathrm{hom}}
\newcommand{\alg}{\mathrm{alg}}
\title{Local Green Forms and Chow Correspondences\\Conditional Reductions and Construction Draft}
\author{}
\date{17 November 2026}
\hypersetup{pdftitle={Local Green Forms and Chow Correspondences},pdfauthor={}}

\begin{document}
\maketitle
\begin{abstract}
We organize two geometric reductions for actual non-Archimedean Green forms. Algebraically trivial classes are treated through additive Poincar\'e families, rational-fibre rigidity and compactness of $p$-adic abelian groups. Homologically trivial classes on abelian varieties are treated through vertical corrections on semistable models; Chow correspondences then transport the resulting numerical certificates. The principal conclusions are stated relative to their actual-form and numerical-transfer hypotheses. The complete available construction arguments are included: model Chern forms, scalar--curvature finite trace, exact-support Green pairs, compact-parameter primitives, abelian numerical descent and the regular-model Chow criterion. Products of curves and abelian varieties occur as consequences of the latter structure.
\end{abstract}

\tableofcontents

\section{The two geometric reductions}\label{sec:construction-introduction}

Fix a finite extension $K/\mathbf Q_p$ and its absolute value. Let $X/K$ be smooth, projective and geometrically integral of dimension $n$. We seek a Green form for a prescribed rational cycle $Z$ of codimension $r$ on the exact open set
\[
 U_Z=X^{\an}\setminus |Z|^{\an},\qquad 1\le r\le n.
\]
Here $|Z|$ is the union of the components with nonzero coefficient. Let $\B^{p,q}$ denote the ordinary actual $\delta$-form sheaf of Mihatsch \cite{M}, and write $D=d'd''$. The desired identity is
\begin{equation}\label{eq:main-green}
 \int_{U_Z}\gamma_Z\wedge D\eta
 =-\int_{Z^{\an}}\eta,
 \qquad \eta\in\B_c^{n-r,n-r}(X^{\an}).
\end{equation}
Its integral on the right retains the rational coefficients and cycle multiplicities. The form on the left belongs to $\B^{r-1,r-1}(U_Z)$, and the identity includes the existence of the displayed pairings.

There are two geometric ways to remove the curvature from a Green pair. Algebraic equivalence supplies a compact parameter group on which the obstruction becomes locally constant. A homologically trivial class on an abelian variety instead supplies a numerical special-fibre correction. The latter is transported through a Chow diagonal factorization. We state these reductions before developing the local constructions they require.

\begin{definition}[Actual obstruction space]\label[definition]{def:construction-obstruction}
Set
\[
 H^{r,r}_{\BC}(X)=
 \frac{\{\omega\in\B^{r,r}(X^{\an}):d'\omega=d''\omega=0\}}
 {D\B^{r-1,r-1}(X^{\an})}.
\]
A prescribed curved Green pair for $Z$ consists of $G_Z\in\B^{r-1,r-1}(U_Z)$ and a global separately closed $\Omega_Z\in\B^{r,r}(X^{\an})$ satisfying
\begin{equation}\label{eq:main-pair}
 \int_{U_Z}G_Z\wedge D\eta+\int_{Z^{\an}}\eta
 =\int_{X^{\an}}\Omega_Z\wedge\eta
 \qquad\left(\eta\in\B_c^{n-r,n-r}(X^{\an})\right).
\end{equation}
\end{definition}

\begin{hypothesis}[Actual curvatures and geometric transfer]\label[hypothesis]{hyp:construction-calculus}
The following properties are named separately, so each conclusion can retain precisely the properties it uses.
\begin{description}
\item[(A1)] A rational-linear, Chow-invariant map
$\Ob^r_{V,K}:\CH^r(V)_{\mathbf Q}\to H^{r,r}_{\BC}(V)$ is computed by actual degree-$r$ Chern characters of virtual bundles on projective flat models. It is functorial for the genuine analytic pullbacks used below. Two model characters with the same generic degree-$r$ Chow class differ by $D$ of a global actual form.
\item[(A2)] Every prescribed rational cycle has a curved Green pair as in \eqref{eq:main-pair}, with $[\Omega_Z]=\Ob^r_{X,K}([Z])$.
\item[(R)] Actual forms on $(X\times S)^{\an}$ have the uniform rational-fibre constancy of \Cref{r12p:rigid-fiber}: for smooth separated $S/K$ and $a_0\in S(K)$, the pulled-back actual form is constant for all $a\in S(K)$ in one ordinary neighbourhood of $a_0$.
\item[(N)] Let $h:V\to X$ be finite dominant between smooth integral projective $K$-varieties. If a projective flat model of $V$ carries $\xi\in K_0(\mathcal V)_{\mathbf Q}$ such that
\[
 \operatorname{ch}_r(\xi|_V)=h^*z,\qquad
 \deg\bigl(\operatorname{ch}_r(\xi_s)\cap[T]\bigr)=0
\]
for every integral vertical $r$-cycle $T$, then $\Ob^r_{X,K}(z)=0$. A variety over a finite larger constant field may be regarded as a $K$-variety; the special-fibre multiplicities, residue degrees and full native degree are retained.
\item[(C)] The kernels of the obstruction maps are preserved by algebraic Chow correspondences of the specified degree, including proper algebraic pushforward. For a finite extension $L/K$, this is interpreted in the native $K$-theory and the algebraic projection $\pi:X_L\to X$ satisfies $\pi_*\pi^*=[L:K]$.
\end{description}
These properties are considered for every smooth projective integral variety occurring in the relevant argument, componentwise after geometric base change when its constant field is larger.
\end{hypothesis}

\begin{theorem}[Algebraic equivalence]\label[theorem]{thm:main-algebraic}
Assume (A1), (A2) and (R). If $[Z]$ is algebraically equivalent to zero over $\overline K$, then $\Ob^r_{X,K}([Z])=0$. Consequently there exists $\gamma_Z\in\B^{r-1,r-1}(U_Z)$ satisfying \eqref{eq:main-green} over the original valued field and against every compact actual complementary test.
\end{theorem}

The proof is completed in \Cref{sec:conditional-algebraic-proof}. The Poincar\'e family converts the cycle into an additive family on a Jacobian. Its algebraic norm gives a family over $K$ with one fibre equal to $[L:K][Z]$. The obstruction-valued map is locally constant by (R), and compactness annihilates it.

For the second reduction, a degree-$a$ correspondence $V\rightsquigarrow W$ means a class in $\CH^{\dim V+a}(V\times W)_{\mathbf Q}$, acting from codimension $r$ to codimension $r+a$.

\begin{definition}[Potential abelian Chow retraction]\label[definition]{def:construction-retraction}
The variety $X$ has a potential abelian Chow retraction if there are a finite extension $L/K$, abelian varieties $A_j/L$, integers $a_j$, and correspondences
\[
 u_j\in\CH^{n+a_j}(X_L\times_LA_j)_{\mathbf Q},\qquad
 v_j\in\CH^{\dim A_j-a_j}(A_j\times_LX_L)_{\mathbf Q}
\]
such that
\begin{equation}\label{eq:construction-retraction}
 \Delta_{X_L}=\sum_jv_j\circ u_j
 \quad\text{in }\CH^n(X_L\times_LX_L)_{\mathbf Q}.
\end{equation}
The zero-dimensional abelian variety $\Spec L$ is allowed.
\end{definition}

\begin{theorem}[Abelian correction and Chow transfer]\label[theorem]{thm:main-abelian}
Assume (A1), (A2), (N), (C), and the semistable abelian model theorem recalled in \Cref{thm:kunnemann-local}. If $X$ has a potential abelian Chow retraction, then
\[
 \Ob^r_{X,K}\bigl(\CH^r(X)_{\mathbf Q,\homol}\bigr)=0.
\]
For $1\le r\le n$, every prescribed cycle in this group therefore has a Green form satisfying \eqref{eq:main-green}. In a fixed codimension, it is enough that the right-hand side of \eqref{eq:construction-retraction} act on the homologically trivial group as a fixed nonzero rational multiple of the identity.
\end{theorem}

Here $\CH^r(X)_{\mathbf Q,\homol}$ is the kernel of the geometric rational $\ell$-adic cycle class for a fixed $\ell\ne p$. The proof, completed in \Cref{sec:conditional-abelian-proof}, first produces a vertically numerical extension on a semistable abelian model. Property (N) converts that extension into zero obstruction. Property (C) then applies the actual Chow identity term by term.

\begin{corollary}[Products of curves and abelian varieties]\label[corollary]{cor:main-products}
Under the hypotheses of \Cref{thm:main-abelian}, let
\[
 X=A\times_KC_1\times_K\cdots\times_KC_s,
\]
where $A/K$ is abelian and the $C_i/K$ are smooth projective geometrically integral curves. Every prescribed geometrically homologically trivial cycle of codimension $1\le r\le\dim X$ satisfies \eqref{eq:main-green} over $K$.
\end{corollary}
\begin{proof}
After a finite extension choose rational points on all the curves. The actual Chow diagonal factorization of \Cref{lem:curve-product-diagonal} factors through the single abelian variety $A\times\prod_i\operatorname{Jac}(C_i)$ with Tate shifts. Apply \Cref{thm:main-abelian}.
\end{proof}

The available construction of correspondence transfer also has a model-sensitive form. \Cref{thm:chow-abelian-criterion} requires a regular integral projective flat target model after finite extension, together with a regular semistable abelian model and the Chow identity. \Cref{lem:product-model} supplies the models for products. This realizes the product argument through concrete numerical certificates, rather than assuming the abstract kernel property (C) for every correspondence.

\subsection{How the local construction is assembled}\label{sec:construction-strategy}

The curvature construction precedes the support construction. Model Chern forms and the numerical detector are developed in \Cref{r10c:section}. The restricted trace in \Cref{r11p:minimal-trace} then handles original scalar--curvature generators; its one-parameter formula recovers the scalar potential without expanding its domain to arbitrary actual forms. In \Cref{r11p:supported-complex}, supported perfect complexes and finite-flat multisections produce the prescribed residue and exact-open Green pair. The construction draft derives (A1) by producing global primitives for generic Chow-zero characters.

These operations are applied in two ways. \Cref{r12p:section} constructs primitives from compact additive families. \Cref{sec:abelian-numerical} constructs them from semistable numerical certificates and traces the original generator expressions along composite finite maps. Finally, \Cref{sec:correspondences} transports certificates by algebraic intersection theory on regular models and proves the product consequences.

\section{Conventions and cited foundations}\label{sec:foundations}

For a projective flat model $\mathcal X/K^\circ$, a model line bundle $\mathcal L$ determines a PL metric on its generic fibre. Its actual curvature $c_1(\overline L)=D(-\log\|e\|)$ has pure tridegree $(0,0,1)$; pure tridegree $(0,0,r)$ has bidegree $(r,r)$. Special-fibre intersections are rational and operational. Native tropical calibrations and multiplicities are retained in integration. Finite extensions carry the extension of the fixed absolute value.

The local arguments use Mihatsch's ordinary sheaf and analytic pullback \cite[Definition 4.2 and Theorem 4.5]{M}, the witnessing-skeleton, restriction and support statements \cite[Lemma 4.4, Corollary 4.9 and Proposition 4.11]{M}, native finite-flat top-degree integration \cite[Lemma 3.15 and Definition 4.10]{M}, and scalar Poincar\'e--Lelong \cite[Theorem 5.5 and Corollary 5.6]{M}. Polyhedral projection and pullback use \cite[Section 2 and Proposition/Definition 4.2]{MTrop} and \cite[Proposition 2.38 and Definition 2.41]{M}. The analytifications in these applications are pure-dimensional and boundaryless.

The model detector uses scalar uniform density of rational model functions on proper analytifications \cite[Proposition 2.15]{GM}. Its measure calculation uses the recursive model intersections and component-mass formula of Chambert-Loir--Ducros \cite[Sections 8.7--8.8 and Theorem 18.5.5]{CLD}. The localized Riemann--Roch calculation is the Cartier case of \cite[Corollary 18.1.2]{Fulton}.

The algebraic arguments use rational Chern characters on regular schemes, vector-bundle resolutions on smooth projective varieties, supported proper and refined Gysin identities, Bertini, miracle flatness and characteristic-zero resolution. The relevant references are \cite{Stacks,StacksFEY,StacksFEW,StacksGysin}. Curve-parameter generation of algebraic equivalence is used as in \cite{ACMV}, and normalized Poincar\'e bundles as in \cite{MilneJV}. The precise semistable component statement needed for abelian varieties is recalled in \Cref{thm:kunnemann-local}.


% Source: research-sources/green-draft/model_chern.tex
\section{Actual Chern forms of model vector bundles}
\label{r10c:section}

Let $K$ be a complete nontrivially discretely valued field, let
$R=K^\circ$, and let $\mathcal X$ be an integral projective flat
$R$-scheme with smooth projective generic fiber $X$ of pure dimension $n$.
All analytic forms below are Mihatsch actual delta-forms on the ordinary
topology of $X^{\rm an}$; their sheaf is denoted $\mathcal B$ and
$D=d'd''$.  In the intrinsic trigrading, a form in
$\mathcal B^{0,0,r}$ has bidegree $(r,r)$.
The special fiber and its components need not be smooth, reduced, or
normal.  All Chow operations are rational and operational unless
otherwise stated.

\subsection{A scalar-density detector for pure delta forms}

\begin{lemma}[Detection with model Hessians]
\label{r10c:detector}
Let $\Omega\in\mathcal B^{0,0,r}(X^{\rm an})$ satisfy
$d'\Omega=d''\Omega=0$, with $0\leq r\leq n$, and put $q=n-r$.  Suppose that for every
tuple of global rational model functions $\phi_1,\ldots,\phi_q$,
\[
                 \Omega\wedge D\phi_1\wedge\cdots\wedge D\phi_q=0
\]
as a signed Radon measure.  Then $\Omega=0$ as an actual form.
For $q=0$, the hypothesis is the vanishing of the measure of $\Omega$.
\end{lemma}

\begin{proof}
An actual top-degree form has a locally finite signed measure: in finitely
many polyhedral charts over a compact set this follows by integration of
its smooth coefficient densities.  A partition of unity gives the same
measure as actual integration.  If that measure is zero, Mihatsch
Proposition 4.11 implies that the top form is zero.  This deals with $q=0$.

Assume that vanishing of the displayed measure has already been proved
with $k$ model-function slots replaced by arbitrary global original-smooth
functions.  Fix those $k$ functions and all but one of the remaining
model functions, and denote their product with $\Omega$ by $T$.
Then $T$ has bidegree $(n-1,n-1)$ and is separately closed.  For every
model function $\phi$ and every global original-smooth scalar $\psi$,
two applications of actual Stokes give
\[
 \int_{X^{\rm an}}\phi\,T\wedge D\psi
     =\int_{X^{\rm an}}\psi\,T\wedge D\phi=0.
\]
For fixed $T,\psi$, the left-hand side tests the single fixed finite
signed measure $T\wedge D\psi$ against $\phi$.  Scalar uniform density
therefore implies $T\wedge D\psi=0$.  Repeating this argument proves
\[
          \Omega\wedge D\chi_1\wedge\cdots\wedge D\chi_q=0
\tag{D}\label{r10c:smooth-hessians}
\]
for arbitrary global original-smooth functions $\chi_i$.
The density step concerns the continuous function \(\phi\) against one fixed signed measure, and consequently requires no convergence of Hessians.

If $\Omega$ is nonzero, choose an $n$-dimensional witnessing skeleton
in a moment chart, using Mihatsch Definition 3.8 and Lemma 4.4.
On an open polyhedron $\sigma$ of dimension $q$ its pure tridegree
means that the nonzero coefficient is scalar.  Refine the chart and
subdivide so that its moment coordinates $u_a=-\log|h_a|$ have
injective tangent map on $\sigma$; this follows from the finite-fiber
property of skeleton tropicalization.  At a point in the relative
interior of $\sigma$, choose a smooth cutoff equal to one nearby and
supported in the moment-chart open set.  It globalizes
\[
                         \chi=\tfrac12\sum_a u_a^2
\]
to an original-smooth scalar on $X^{\rm an}$.  Near the chosen point,
$D\chi=\sum_a d'u_a\wedge d''u_a$ is positive definite on the tangent
space of $\sigma$.  Thus $(D\chi)^q$ is a nonzero volume form there.
Coefficientwise wedge with smooth forms now shows that
$\Omega\wedge(D\chi)^q$ is a nonzero actual top form.  This contradicts
\eqref{r10c:smooth-hessians} and Proposition 4.11.
\end{proof}

\subsection{From operational numerical triviality to actual zero curvature}

For clarity, the CLD top measure agrees with the measure of the actual
wedge used above.  Locally let $u$ be a PWS model metric weight and let
$T$ be a product of previously treated actual model curvatures.  Such a
current is locally of order zero by its polyhedral description, so
coefficientwise multiplication by the continuous function $u$ is defined
and $[uT]=u[T]$. This equality is tested against original-smooth forms
and follows directly from the same order-zero polyhedral integrals
on both sides. Separate closedness and Mihatsch Proposition 4.11 give
\[
                  [Du\wedge T]=D[uT]=D\bigl(u[T]\bigr).
\]
Induction first identifies the actual wedge with the CLD recursive
product for semipositive model factors. For signed model metrics,
express them as differences of psh model
metrics and use multilinearity.  This establishes the comparison over
$K$ itself, before any scalar extension.

\begin{theorem}[Numerical model curvature vanishing]
\label{r10c:numerical-zero}
Let $P$ be a homogeneous polynomial of degree $r$ in the first Chern
classes of model line bundles $\mathcal L_1,\ldots,\mathcal L_t$ on
$\mathcal X$.  Suppose that
\[
       \xi=P\bigl(c_1(\mathcal L_{1,s}),\ldots,
                    c_1(\mathcal L_{t,s})\bigr)\in A^r(\mathcal X_s)_{\mathbb Q}
\]
has degree zero against every integral $r$-dimensional closed subvariety
of $\mathcal X_s$.  Then
\[
        P\bigl(c_1(\overline L_1),\ldots,c_1(\overline L_t)\bigr)=0
                 \quad\hbox{in }\mathcal B^{0,0,r}(X^{\rm an}).
\]
\end{theorem}

\begin{proof}
Let $\Omega$ be the displayed actual model polynomial.  Given $q=n-r$
model functions, choose a common proper model refinement
$f:\mathcal Y\to\mathcal X$ on which they are represented, after
clearing denominators, by vertical Cartier divisors $D_1,\ldots,D_q$.
One may work with formal completions; algebraic models and projective
common refinements suffice here.  The induced metrics of the original
bundles are unchanged by pullback.

For a point $x$ reducing to the generic point of a special component
$V\subset\mathcal Y_s$, the top wedge
$\Omega\wedge\prod_jD\phi_j$ has mass equal to the native CLD weight
$w_{\mathcal Y}(x)$ times
\[
 \deg\left(f_s^*\xi\cap
       \prod_{j=1}^{q}c_1(\mathcal O(D_j)|_V)\cap[V]\right).
\]
The proper operational projection formula identifies this degree with
\[
 \deg\left(\xi\cap f_{s*}
       \left(\prod_{j=1}^{q}c_1(\mathcal O(D_j)|_V)\cap[V]\right)\right)=0.
\]
The pushed-forward class has dimension $r$.  Terms whose images have
smaller dimension have zero proper pushforward; all remaining terms are
annihilated by the numerical hypothesis.  Operational pullback of a
Cartier polynomial is available for arbitrary $f_s$: neither flatness of
$f_s$ nor smoothness of $V$ has been used.  The CLD support and mass
formula now gives a zero measure, with all native multiplicities and
defect factors retained.  Lemma~\ref{r10c:detector} concludes the proof.
\end{proof}

\subsection{Descent along locally split morphisms}

\begin{lemma}[Descent for a locally split morphism]
\label{r10c:split-descent}
Let $a:Y\to X$ be a morphism of pure-dimensional boundaryless analytic
spaces.  Suppose $X$ has an ordinary open cover $(U_i)$ and morphisms
$s_i:U_i\to Y$ with $a\circ s_i=\operatorname{id}_{U_i}$.
For $\Theta\in\mathcal B(Y)$ the following are equivalent:
\begin{enumerate}
\item There is an actual form $\theta$ on $X$ with $a^*\theta=\Theta$.
\item On $Y\times_XY$, $p_1^*\Theta=p_2^*\Theta$.
\end{enumerate}
When the fiber product is pure-dimensional and the displayed pullbacks
are defined, the descended form is unique and is given by
$\theta|_{U_i}=s_i^*\Theta$.
\end{lemma}

\begin{proof}
The forward implication is functoriality.  Conversely pull the equality
back along $(s_i,s_j):U_i\cap U_j\to Y\times_XY$ to obtain equality
of the two local forms.  They glue by the sheaf property.  Pulling the
same equality along
$(\operatorname{id}_Y,s_i\circ a):a^{-1}(U_i)\to Y\times_XY$ proves
$\Theta=a^*s_i^*\Theta$ on $a^{-1}(U_i)$, hence $\Theta=a^*\theta$.
The local sections also prove uniqueness and injectivity of $a^*$ on
actual forms.  
\end{proof}

\begin{theorem}[Actual model Chern forms by local flag sections]
\label{r10c:chern}
Let $\mathcal E_1,\ldots,\mathcal E_s$ be vector bundles on
$\mathcal X$ and $P$ a homogeneous rational polynomial of degree $r$
in their Chern classes.  There is a unique actual form
\[
              c_{\mathcal X}(P)\in\mathcal B^{0,0,r}(X^{\rm an})
\]
whose pullback to the simultaneous full flag bundle equals the
corresponding symmetric polynomial in the model first Chern forms of
the tautological line quotients.  The form is separately closed,
independent of the chosen splitting presentation, and functorial under
morphisms of projective flat models with smooth projective generic
fibers.  Universal operational identities, in particular Whitney and
the additivity of the Chern character in exact sequences of model
bundles, hold as strict identities of actual forms.

If the special-fiber operational class $P(c(\mathcal E_{j,s}))$ is
numerically zero against every $r$-cycle, then $c_{\mathcal X}(P)=0$.
\end{theorem}

\begin{proof}
Put $a:\mathcal F=\prod_{\mathcal X}\operatorname{Flag}(\mathcal E_j)
\to\mathcal X$ and denote its generic fiber by $a:F\to X$.
Let $\Omega_F$ be the symmetric model-root polynomial.  On
$\mathcal F\times_{\mathcal X}\mathcal F$, the two root polynomials
represent the same operational Chow class by the two tautological
filtrations and Whitney's identity.  Their difference is therefore
operationally zero on the special fiber, before taking a degree or
capping with the fundamental class.  Theorem~\ref{r10c:numerical-zero},
applied to the smooth projective generic fiber $F\times_XF$, gives
\[
                         p_1^*\Omega_F=p_2^*\Omega_F
\]
as actual forms.

Simultaneously trivialize the finitely many generic bundles $E_j$ on
a Zariski open cover of $X$.  Their standard full flags give algebraic
sections $s_i:U_i\to F$, hence analytic sections on ordinary opens
$U_i^{\rm an}$.  These sections need not extend to all of the model
and need not preserve a model trivialization.  This is harmless:
they pull back an already defined global actual form, and Mihatsch
Theorem 4.5 applies to arbitrary analytic morphisms.  Applying
Lemma~\ref{r10c:split-descent} defines
$c_{\mathcal X}(P)|_{U_i^{\rm an}}=s_i^*\Omega_F$ and proves its
characterizing pullback equality and uniqueness.  Pullback preserves
the intrinsic trigrading and the differential operators, which proves
purity and separate closedness.

For any operational identity among finitely many model bundles, use
their simultaneous flag bundle.  The difference of the root polynomials
is operationally zero, so the numerical detector kills it upstairs.
Faithfulness supplied by local sections gives the identity downstairs.
The same argument on the product of two flag presentations proves
independence of presentation.  Functoriality follows by pulling the
model flags back and then using their characteristic equality and the
faithfulness of the new flag pullback.

Finally suppose $\xi=P(c(\mathcal E_{j,s}))$ is numerically zero.
For every integral $r$-cycle $W$ on $\mathcal F_s$, the operational
projection formula gives
\[
                 \deg(a_s^*\xi\cap[W])=
                 \deg(\xi\cap a_{s*}[W])=0.
\]
If the image drops dimension then $a_{s*}[W]=0$; otherwise the numerical
hypothesis applies.  Thus the root polynomial is numerically zero on
the whole flag special fiber.  Theorem~\ref{r10c:numerical-zero}
gives $\Omega_F=0$, and local sections give $c_{\mathcal X}(P)=0$.
Although $a_s$ happens to be smooth and flat here, the projection
argument is the operational one and does not use a smooth special fiber.
\end{proof}

The construction gives a homomorphism
\[
 K_0(\mathcal X)_{\mathbb Q}\longrightarrow
      \bigoplus_{r=0}^{n}\mathcal B^{0,0,r}(X^{\rm an}),
                  \qquad [\mathcal E]\longmapsto
                       \sum_r c_{\mathcal X}(\operatorname{ch}_r\mathcal E).
\]
Each degree factors through its special-fiber operational numerical
class.  This statement concerns classes of model bundles.  The dependence on the model is controlled by the numerical class of its special fiber.



% Source: research-sources/green-draft/finite_trace.tex
\section{Finite traces of scalar--curvature forms}
\label{r11p:minimal-trace}

The trace will be constructed for the scalar--curvature module generated by products
\[
  u_0\,Du_1\cdots Du_k,\qquad D=d'd'',
\tag{MT1}\label{r11p:module-generators}
\]
where the $u_i$ are piecewise linear analytic potentials.

Throughout this section $K$ is a complete nontrivially valued field of
characteristic zero, and all analytic spaces are the analytifications of
normal integral separated $K$-varieties, with their ordinary topology.
All forms are Mihatsch actual delta-forms.  A PL scalar means a function
locally presented by a finite piecewise affine function of logarithms of
analytic units.  Its $D$ is a pure form of tridegree $(0,0,1)$.
Let $\mathscr M^k$ be the actual image sheaf of finite sums of
\eqref{r11p:module-generators}.  Relations in $\mathscr M^k$ are actual-form
relations; no chosen presentation is part of a section.

\subsection{Native degree under finite maps}

For a finite dominant map $q:Y\to X$ of generic degree $d$ and any compact
actual top form $\tau$ on $X$, these inputs give
\[
                            \int_Yq^*\tau=d\int_X\tau.
\tag{MT2}\label{r11p:native-degree}
\]
Indeed remove a proper closed subset outside which $q$ is finite flat of
constant degree $d$.  The support of $\tau$ avoids this subset by
Corollary 4.9; so does the support of its pullback above it.
Lemma 3.15 and the native skeletal integration formula apply on the
remaining open.    It also proves that finite dominant pullback on
actual forms is faithful: pair a putatively zero pullback with pulled-back
original-smooth complementary tests, use \eqref{r11p:native-degree}, and
apply Proposition 4.11.

\subsection{Descent of scalar PL functions through finite quotients}

\begin{lemma}[Ordinary local scalar PL descent]
\label{r11p:pl-descent}
Let $q:Y\to X$ be finite Galois between normal integral varieties, with
group $G$.  A $G$-invariant source PL scalar descends to a PL scalar on
$X$.  If finitely many invariant potentials on
$Y^{\rm an}\times(\mathbf G_m)^{\rm an}$ are jointly PL in fixed source
unit logarithms and the parameter $t=v(z)$, their descended potentials
are locally jointly PL in a fixed finite tuple of base-unit logarithms
and $t$.  The base tuple and base shrinking are independent of $t$.
\end{lemma}

\begin{proof}
Work around a whole finite analytic fiber.  Choose disjoint ordinary
neighborhoods of its points on which the finitely many source potentials
have moment presentations.  Properness permits shrinking the base so
that its full inverse image lies in their union.  The corresponding
source pieces are then clopen.  Analytic units can be pasted on these
pieces.  Add constant units with distinct valuations as markers; after
shrinking compact neighborhoods, their moment images lie in separated
compact polyhedral boxes.  Finite PL extension on these boxes
combines the presentations.  A smooth cutoff of a PL potential is
neither required nor used.

Include the full $G$-orbit of the resulting source-unit tuple.  Its
logarithms $\ell\in\mathbf R^s$ then carry a coordinate-permutation
$G$-action; add a constant unit if necessary so that $s\geq1$.
If an ambient PL presentation $\varphi(\ell)$ is
invariant only after restriction to the actual source moment image,
replace it by $|G|^{-1}\sum_{g\in G}\varphi(g\ell)$.  This is an
ambient invariant PL function with the same source restriction.
For ordered orbit potentials the sorting itself already gives this
ambient invariance.  The same operation leaves an independent
parameter fixed.  Put $d=|G|$ and $N=d(s-1)+1$.  For distinct integers $a_j$
take the Vandermonde directions
\[
                     m_j=(1,a_j,\ldots,a_j^{s-1})\in\mathbf Z^s.
\]
Define $J(\ell)$ by recording, for each $j$, the ordered multiset
$\{\langle m_j,g\ell\rangle:g\in G\}$.  It separates $G$-orbits.
In fact if $J(x)=J(y)$, for each $j$ one can choose $g_j$ with
$\langle m_j,x\rangle=\langle m_j,g_jy\rangle$.  One $g$ occurs for
at least $s$ indices.  The resulting polynomial of degree at most $s-1$
has $s$ roots, hence $x=gy$.

The map $J$ is proper: any $s$ independent Vandermonde directions bound
$\ell$ when all the displayed coordinates are bounded.  It is finite
piecewise linear.  On a common polyhedral subdivision its linearity
cells map affinely and injectively to their images; a noninjective
tangent direction would produce an infinite fiber.  Every invariant
PL scalar therefore factors as a PL scalar through $J$, after common
image subdivision and PL ambient extension.  The identical argument
applies to $J\times\operatorname{id}_{\mathbf R}$ and jointly PL
potentials with the fixed parameter $t$.

It remains to realize $J$ downstairs without an illicit nonlinear
pullback of forms.  For each source monomial $F_j=f^{m_j}$ take its monic
orbit polynomial
\[
                          \prod_{g\in G}(T-gF_j).
\]
Its coefficients are invariant analytic functions, hence analytic
functions on the normal finite quotient.  Algebraically the quotient
is $X$; analytically this follows on affinoids from finite coherent
base change and the Reynolds idempotent in characteristic zero.
With the convention that the coefficient of $T^i$ gives the point
$(i,v(a_i))$, the ordered root valuations are the negatives of the
lower Newton polygon's slopes, with multiplicity.  They are therefore
PL in the valuations of the coefficients that contribute to that
polygon.

At a fixed base point retain the nonzero coefficients and shrink so
that they are units.  The constant coefficient $a_0$ is a unit because
it is a product of source units; the leading coefficient is one.
Choose a uniform upper bound $B$ for $v(a_0)$ on the shrunken base.
Every coefficient vanishing at the fixed point can, by continuity and
a further common shrinking, be made to have valuation greater than
$\max(B,0)+1$.  Its point is then above the endpoint chord of the
Newton polygon and cannot affect the lower hull.  Such coefficients
can be omitted uniformly.
Thus $J$ is PL in finitely many base-unit logarithms.  These
coefficients involve source units alone, so their choice and shrinking
are independent of $t$.  Substitution is only a substitution into a
scalar PL function.  Taking its actual derivatives happens afterwards;
no nonlinear PL pullback of a general delta-form has been postulated.
\end{proof}

\subsection{Sorting gives a pure-curvature trace}

\begin{lemma}[Max/min identities]
\label{r11p:sorting}
For source PL scalars $u,v$, put $M=\max(u,v)$ and $m=\min(u,v)$.
For every $k\geq0$ there are actual identities
\[
 (DM)^k+(Dm)^k=(Du)^k+(Dv)^k,
\tag{MT3}\label{r11p:sorting-powers}
\]
\[
 M(DM)^k+m(Dm)^k=u(Du)^k+v(Dv)^k.
\tag{MT4}\label{r11p:sorting-energies}
\]
\end{lemma}

\begin{proof}
Write $h=u-v$, $h_+=\max(h,0)$, $h_-=\min(h,0)$, and
$A=Dh_+$, $B=Dh_-$.  Coefficientwise multiplication on support
polyhedra gives $h_+B=h_-A=0$: on one open side the appropriate
function is zero, and on the common wall its continuous scalar
coefficient is zero.  Applying $D$ and using separate closedness gives
$AB=0$.  Since
\[
 M=v+h_+,\quad m=v+h_-,\quad u=v+h_++h_-,
\]
expansion proves \eqref{r11p:sorting-powers}.  The same expansion,
now using both $AB=0$ and the two scalar-weighted zero identities,
proves \eqref{r11p:sorting-energies}.
\end{proof}

The second identity is useful but has a definite limitation:
polarizing $u(Du)^k$ recovers symmetrized energies, not an arbitrary
ordered scalar-curvature monomial $u_0Du_1\cdots Du_k$.
The parameter argument below supplies precisely that missing step.

\begin{lemma}[Trace of pure PL curvature polynomials]
\label{r11p:pure-trace}
For finite dominant $q:Y\to X$ between normal integral varieties, a
finite sum $C$ of products $Du_1\cdots Du_k$ has an intrinsic actual
trace $\operatorname{Tr}^{\rm curv}_q(C)$, characterized by
\[
 \int_X\operatorname{Tr}^{\rm curv}_q(C)\wedge\eta
                         =\int_YC\wedge q^*\eta
\tag{MT5}\label{r11p:pure-pairing}
\]
for all compact complementary actual tests.  With a parameter, jointly
PL source potentials in fixed source logarithms and $t$ yield a trace
with a product-moment presentation using fixed base logarithms and $t$.
\end{lemma}

\begin{proof}
First let $q$ be Galois.  Sort the orbit of a scalar $u$ by a finite
sorting network.  Repeated use of \eqref{r11p:sorting-powers} gives
\[
                     \sum_{g\in G}(Dg^*u)^k
                            =\sum_{j=1}^d(Du_{(j)})^k.
\]
The ordered potentials $u_{(j)}$ are invariant and descend by
Lemma~\ref{r11p:pl-descent}.  The right side is consequently the
pullback of a downstairs actual form.  Define it to be the trace of
$(Du)^k$.  Polarization,
\[
 Du_1\cdots Du_k
   =\frac1{k!}\sum_{I\subset\{1,\ldots,k\}}
           (-1)^{k-|I|}\left(D\sum_{i\in I}u_i\right)^k,
\]
does the same for all mixed products.  For $k=0$ set
$\operatorname{Tr}^{\rm curv}_q(1)=d$.

This construction satisfies the actual orbit identity
\[
                 q^*\operatorname{Tr}^{\rm curv}_q(C)
                                      =\sum_{g\in G}g^*C.
\]
Wedge with $q^*\eta$, integrate, and use
\eqref{r11p:native-degree} and invariance of native integration under
automorphisms.  Division by $d$ proves \eqref{r11p:pure-pairing}.
Finite pullback faithfulness proves independence of every actual
relation.  Current injectivity gives gluing on ordinary overlaps.
The pairing permits arbitrary actual downstairs multipliers; they
have not been treated as arbitrary source coefficients.

For a general separable finite map take a normal Galois closure
$\widetilde Y\to X$ factoring through $r:\widetilde Y\to Y$, and divide
the Galois trace of $r^*C$ by $e=\deg r$.  Applied to the top form
$Cq^*\eta$, \eqref{r11p:native-degree} proves the required pairing and
independence of the closure.  Characteristic zero covers every map
used here.  All spaces and parameter tori remain over the original
$K$; native generic degrees are not normalized away.

For the last assertion use the parameter version of
Lemma~\ref{r11p:pl-descent} at the sorting step.  The descended
potentials have fixed base moments, so their Hessians and products
have the required product-moment presentation.
\end{proof}

\subsection{Integration along a parameter torus}

\paragraph{The boundary issue.}
Compact analytic support does not imply that an ambient moment
coefficient avoids the tropicalized analytic boundary of an affinoid
chart.  We do not use the single-chart integral formula with that
unverified hypothesis.  Instead we first prove a pushforward identity
for polyhedral currents on the product skeleton, retaining the boundary
exclusion of Proposition 2.38 at every chart calculation.  The
intrinsic skeletal definition of integration then gives the analytic
pairing.  In this subsection $f^\star$ denotes tropical-space pullback
in that proposition, distinguished from the notation $t_c^*$ for the
resulting actual analytic form.

\begin{lemma}[The coordinate product skeleton, including its structures]
\label{r11p:coordinate-product-skeleton}
Let $V$ be a boundaryless ordinary analytic open of pure dimension $n$,
let $b$ be a tuple of analytic units on $V$, and put $S=\Sigma(V,b)$.
The relative Gauss section identifies
\[
 \begin{split}
 \iota_b:S\times\mathbf R&\ \xrightarrow{\ \sim\ }\
       \Sigma(V\times(\mathbf G_m)^{\rm an},(b,z)),\\
                  (y,t)&\longmapsto \nu_t(y),
 \end{split}
\tag{PS1}\label{r11p:product-skeleton-identity}
\]
where $\nu_t(y)$ is the Gauss point of radius $e^{-t}$ in the torus
over $\mathcal H(y)$.  This is an isomorphism of weighted pwl spaces
and of tropical spaces.  Its weight is the product of the native
weight on $S$ and the standard weight on $\mathbf R$, and its linear
function sheaf is generated by the base linear functions and $t$.
The identification is compatible with every ordinary base shrinking
and every additional base moment refinement.
\end{lemma}

\begin{proof}
Let $x$ lie above $y$ in the analytic product.  The graded field
generated by $b$ and $z$ has transcendence degree $n+1$ precisely
when the graded field generated by $b$ at $y$ has transcendence
degree $n$ and $\widetilde z$ is transcendental over
$\widetilde{\mathcal H(y)}$.  Indeed adjoining $z$ increases that
degree by at most one; in the degree-$n$ case
$\widetilde{\mathcal H(y)}$ is algebraic over the graded field
generated by $b$.  This is the graded-transcendence calculation in
(3.27) of the cited proof, here with the single coordinate $z$.

The condition on $\widetilde z$ characterizes the Gauss valuation
over $\mathcal H(y)$: no cancellation among leading monomials is
possible, and for every Laurent polynomial one has
\[
 \left|\sum_i a_i z^i\right|_{\nu_t(y)}
                         =\max_i |a_i|e^{-it}.
\]
Conversely these Gauss valuations have the required graded
transcendence degree.  Thus Definition 3.8 and the absence of analytic
boundary give the claimed bijection.  No valued extension of $K$,
or rationality of the radius in its value group, is needed.

The inverse is $(p,v(z))$ on the displayed skeleton.  It is pwl:
a base skeleton chart is described by an affinoid domain and
valuations of its analytic units; their pullbacks are valuations of
the same units on the product, and $v(z)$ is the extra coordinate.
The corresponding closed chart domains pull back to closed pwl
domains.  This inverse is also proper.  A compact subset of
$S\times\mathbf R$ is contained in a finite union of compact base
analytic domains times a bounded interval.  Its inverse image is
closed in the corresponding finite union of products with compact
analytic annuli, and is therefore compact.  It follows that the
continuous bijective inverse is a homeomorphism.  Subdivision into
local polyhedral charts makes its injective affine pieces invertible
onto their images, so both directions are pwl.  This is the
coordinate specialization of the topology and atlas argument
following (3.31), now without its changes of field or center.

For the weights, take a top-dimensional base cell and a local toric
projection used in Definition 3.8(2) to compute its native degree
$e$.  Its product with the identity on $\mathbf G_m$ is finite flat
of the same degree $e$.  Hence its induced weight is the base weight
times the standard parameter weight.  This calculation retains
residue-field degrees and is exactly the last weight calculation
in the proof of Proposition 3.21.  It is local on cells and hence
determines the entire product weight.

Under the bijection the moments are $(t_b,t)$.
Definition 3.8(3) therefore gives precisely the product linear
function sheaf, not merely a set-theoretic identification.
For additional base moments $b'$, the same Gauss section applies.
The inclusions of base and product skeletons agree under it, and
their weights agree by Remark 3.9.  The new linear functions are the
new base ones together with the same $t$.  This proves the stated
compatibility for all refinements and ordinary restrictions.
\end{proof}

\begin{lemma}[Tropical projection with boundary control]
\label{r11p:tropical-product-push}
Let $S$ be a pure-dimensional tropical space and
$f:S\to\mathbf R^N$ a linear map.  Give $S\times\mathbf R$ the
product weight and linear functions.  Write $p_S$ for its projection
and $\pi:\mathbf R^N\times\mathbf R\to\mathbf R^N$ for the coordinate
projection.  If $\Gamma$ is a Euclidean delta-form with support in
a fixed compact set of parameter values, then
\[
 (p_S)_*\bigl((f,\operatorname{id})^\star\Gamma\bigr)
                                      =f^\star(\pi_*\Gamma)
\tag{PS2}\label{r11p:tropical-product-push-identity}
\]
as polyhedral currents on $S$.  The left side uses only the
proper-on-support pwl-current pushforward of Proposition 2.17.
\end{lemma}

\begin{proof}
The product is a tropical space: on a product polyhedral subdivision,
the vertical balancing conditions are those of $S$, and horizontal
internal facets cancel in pairs with the standard parameter weight.
The support of
$B=(f,\operatorname{id})^\star\Gamma$ is contained in
$S\times J_0$ for a compact interval $J_0$.  Thus it is proper over
$S$, and Proposition 2.17 supplies the polyhedral current
$P=(p_S)_*B$ before any assertion that it is a delta-form.

Work in a polyhedral neighborhood carrying a finite-fiber linear
map $h$ to $\mathbf R^m$, as supplied by the axioms of a sheaf of
linear functions.
Append $f$ and put $g=(f,h)$.  Then $g$ has finite fibers and
$f=A\circ g$, where $A$ is a literal coordinate projection.
Let $H$ be any compact polyhedral subset of the interior of this
neighborhood, and put
\[
 E_H=g(\partial_S H),\qquad
 W_H=\mathbf R^{N+m}\setminus E_H,\qquad
 H_g^\circ=H\setminus g^{-1}(E_H).
\]
The restricted $g:H_g^\circ\to W_H$ is proper.  Write
$C_H=g_*[H_{\dim S},\mu_S]|_{W_H}$.
It is a tropical cycle on $W_H$ by Corollary 2.35.

Choose an interval $J=[a,b]$ containing $J_0$ in its interior.
The boundary of $H\times J$ is contained in
\[
        (\partial_S H\times J)\ \cup\ (H\times\{a,b\}).
\]
Apply the characterizing identity (2.40) of Proposition 2.38 to
this compact product and the refinement $g\times\operatorname{id}$.
On the boundary-free region $W_H\times(a,b)$ it gives
\[
 (g\times\operatorname{id})_*
       \bigl(B|_{H_g^\circ\times(a,b)}\bigr)
       =(C_H\times\mathbf R)\wedge
                             (A\times\operatorname{id})^*\Gamma.
\tag{PS3}\label{r11p:boundary-chart-product}
\]
Both sides vanish near the artificial parameter endpoints because
$\Gamma$ does.  They may therefore be viewed on $W_H\times\mathbf R$,
but the exclusion of $E_H$ in the base remains in force.  There
is no assertion of this chart identity across $g(\partial_S H)$.

Let $\pi_g:W_H\times\mathbf R\to W_H$ be projection.  Functoriality
of pwl-current pushforward and the Euclidean projection formula
give, entirely on $W_H$,
\[
\begin{split}
  g_*\bigl(P|_{H_g^\circ}\bigr)
   &=\pi_{g,*}\!\left((C_H\times\mathbf R)\wedge
                          (A\times\operatorname{id})^*\Gamma\right)\\
   &=C_H\wedge
                  \pi_{g,*}(A\times\operatorname{id})^*\Gamma\\
   &=C_H\wedge A^*(\pi_*\Gamma).
\end{split}
\tag{PS4}\label{r11p:boundary-chart-push}
\]
The last step needs only coordinate Fubini.  For a coordinate
projection $A$, its Euclidean delta-pullback adds free standard
coordinates.  On each weighted carrier of $\Gamma$, fiber
integration in $t$ before or after adding those coordinates
has the same fibers, coefficient integral and quotient weights.
Thus
$\pi_{g,*}(A\times\operatorname{id})^*\Gamma=A^*(\pi_*\Gamma)$.
This carrierwise argument applies to an arbitrary Euclidean
delta-form, with compact parameter support; it does not invoke a
general affine base-change lemma.

The last expression in \eqref{r11p:boundary-chart-push} is exactly
the characterizing expression (2.40) for
$f^\star(\pi_*\Gamma)$.  Varying the compact polyhedral set $H$
inside the fixed finite-fiber $g$-chart, Lemma 2.37 shows that
$P=f^\star(\pi_*\Gamma)$ on that chart.  These charts cover $S$,
which proves \eqref{r11p:tropical-product-push-identity}.
In particular the polyhedral pushforward has the asserted
delta-form representative; this is a conclusion of the calculation.
\end{proof}

\begin{lemma}[Integration on the product-moment subspace]
\label{r11p:product-integration}
Let $p:U\times(\mathbf G_m)^{\rm an}\to U$ be projection.  Suppose
\[
                   \beta=t_{(c,z)}^*\Gamma
\]
for a fixed base analytic-unit tuple $c$, $t=v(z)$, and a Euclidean
delta-form $\Gamma(u,t)$ supported in a fixed compact set of $t$.
With $\pi(u,t)=u$, define
\[
                       I_p(\beta)=t_c^*(\pi_*\Gamma).
\tag{MT6}\label{r11p:product-definition}
\]
This is intrinsic among product-moment presentations and satisfies
\[
             \int_UI_p(\beta)\wedge\eta
                    =\int_{U\times\mathbf G_m}\beta\wedge p^*\eta
\tag{MT7}\label{r11p:product-pairing}
\]
for every compact actual complementary test.  It commutes with
$d',d''$ on these presentations.
\end{lemma}

\begin{proof}
\textit{1. A skeletal identity before integration.}
Euclidean pushforward supplies the actual candidate
$\theta=t_c^*(\pi_*\Gamma)$.
On an arbitrary ordinary $V\subseteq U$, refine the base moments
by tuple concatenation, $b=(c,h)$, and put $S=\Sigma(V,b)$.
Under \eqref{r11p:product-skeleton-identity}, the restriction of
$\beta$ to the product skeleton is
$(t_c|_S,\operatorname{id})^\star\Gamma$.
Lemma~\ref{r11p:tropical-product-push} therefore gives
\[
       \theta|_S
        =(p_S)_*\bigl(\beta|_{S\times\mathbf R}\bigr).
\tag{PS5}\label{r11p:actual-product-skeletal}
\]
This is an equality of intrinsic polyhedral currents on $S$ for
every base refinement.  Its proof used the boundary-excluded chart
identities \eqref{r11p:boundary-chart-product}--\eqref{r11p:boundary-chart-push},
not an integral over a global tropicalization of $V$.

If two product-moment presentations define the same actual $\beta$,
concatenate both base tuples and any further test tuple.  Their
restrictions to every resulting product skeleton agree by Definition
4.2.  Equation \eqref{r11p:actual-product-skeletal} makes their
candidate outputs agree on every corresponding base skeleton.
Lemma 4.4, using all refinements, gives actual equality.  This proves
independence of presentations and their relations, and also their
compatibility on ordinary overlaps, before any current pairing.

\textit{2. Products with arbitrary actual base forms.}
On an open $V$ where an actual form $\eta$ has a moment presentation,
take $b$ to be the concatenation of its moments and $c$, and write
$\eta=t_b^*\eta_b$.  Let $A$ be the coordinate projection selecting
$c$.  The product $\beta\wedge p^*\eta$ has Euclidean coefficient
\[
           (A\times\operatorname{id})^*\Gamma\wedge\pi_b^*\eta_b
\]
in the fixed moments $(b,z)$, with the same compact parameter support.
Euclidean projection and coordinate Fubini give
\[
 \pi_{b,*}\!\left(
      (A\times\operatorname{id})^*\Gamma\wedge\pi_b^*\eta_b\right)
                         = A^*(\pi_*\Gamma)\wedge\eta_b.
\]
Apply \eqref{r11p:actual-product-skeletal} to this coefficient.
For $S=\Sigma(V,b)$ it gives
\[
 (p_S)_*
     \left((\beta\wedge p^*\eta)|_{S\times\mathbf R}\right)
                         =(\theta\wedge\eta)|_S.
\tag{PS6}\label{r11p:product-test-skeletal}
\]
The coefficient $\eta_b$ here is an arbitrary Euclidean delta-form.
No restriction to smooth tests or to products of model curvatures
has occurred.

\textit{3. Intrinsic skeletal integration and ordinary partitions.}
Now let $\eta$ be compact and of complementary bidegree.  Cover its
support by finitely many ordinary opens $V_i$ on which it has the
presentations used in the preceding step, and choose an
original-smooth partition $(\lambda_i)$ subordinate to this cover
on a neighborhood of the support.  Put
$S_i=\Sigma(V_i,b_i)$.
On $V_i$ the unweighted form $\theta\wedge\eta$ is top-degree
and is presented in $b_i$; on $V_i\times\mathbf G_m$ the
unweighted form $\beta\wedge p^*\eta$ is top-degree and is
presented in $(b_i,z)$.  Corollary 4.9(2) puts their supports in
$S_i$ and $S_i\times\mathbf R$, respectively.

Multiplication by $\lambda_i$ or $p^*\lambda_i$ cannot enlarge
these supports.  Their paired supports are compact: the relevant
base support is compact inside $V_i$, and the parameter values lie
in a fixed compact interval, hence in a compact analytic annulus.
Finitely many affinoid neighborhoods of the base support give
compact total paired support.  This compactness justifies the
integrals; it is not used to exclude any tropicalized base boundary.

The restriction $\lambda_i|_{S_i}$ need only be PWS.  Remark 4.3(1)
makes its multiplication coefficientwise, and Proposition 2.17
gives the coefficientwise projection formula
\[
 (p_{S_i})_*\!\left(p_{S_i}^*(\lambda_i|_{S_i})\,T\right)
                  =(\lambda_i|_{S_i})(p_{S_i})_*T
\]
for the relevant polyhedral current $T$.
Use Definition 4.10(1) to evaluate the analytic top integrals on
these containing skeletons, followed by (2.17) and the integrable
current pushforward formula in Definition 4.10(4).  With
\eqref{r11p:product-test-skeletal}, this gives
\[
\begin{split}
 \int_{U\times\mathbf G_m}p^*\lambda_i\,\beta\wedge p^*\eta
 &=\int_{S_i\times\mathbf R}
      p_{S_i}^*(\lambda_i|_{S_i})
          (\beta\wedge p^*\eta)|_{S_i\times\mathbf R}\\
 &=\int_{S_i}(\lambda_i|_{S_i})(\theta\wedge\eta)|_{S_i}\\
 &=\int_U\lambda_i\,\theta\wedge\eta.
\end{split}
\]
Definition 4.10(2) permits summing this identity over the ordinary
partition.  It proves \eqref{r11p:product-pairing} for every
compact actual test.  The partition is not required to have a
single Euclidean presentation in the preassigned tuple $b_i$.

Finally, Euclidean proper-on-support pushforward commutes with
$d',d''$, and differentiation retains the compact parameter support.
Applying this to \eqref{r11p:product-definition}, whose actual
independence is already proved, gives the asserted differential
compatibility.  Equivalently it follows from the linear-map
current identity of Definition 2.29(3) in
\eqref{r11p:actual-product-skeletal}.
\end{proof}

\begin{lemma}[One-parameter evaluation of a nonlinear PL scalar]
\label{r11p:one-parameter}
Let $u$ be a bounded PL scalar on an ordinary open $V$.  Choose
$\chi\in C_c^\infty(\mathbf R)$ equal to one near its range.  On
$V\times(\mathbf G_m)^{\rm an}$ put
\[
                         H=\max(u-t,0).
\]
For every actual form $C$ on $V$ with a common moment presentation,
\[
                  I_p\bigl(\chi(t)t\,DH\wedge p^*C\bigr)=uC.
\tag{MT8}\label{r11p:graph-evaluation}
\]
Only one parameter occurs, independently of the number of source
moments needed to present $u$.
\end{lemma}

\begin{proof}
Write $u=\varphi(\ell)$ for a finite Euclidean PL scalar.  Work first
on its Euclidean moment domain with $H(x,t)=\max(\varphi(x)-t,0)$.
On each cell where $\varphi$ is affine, the graph face $t=\varphi(x)$
in $DH$ has projection weight one.  The coefficient of $t$ in the
defining affine difference is $-1$, so projection of its calibrated
divisor measure is the standard base measure.  Weighting by $\chi(t)t$
gives $\chi(\varphi(x))\varphi(x)$.

There are also vertical faces above walls of $\varphi$, in the region
$t<\varphi(x)$.  It would be false to identify $DH$ with the graph
alone.  On every such vertical face, however, projection drops the
dimension of its pure carrier; its pure Euclidean pushforward is
zero, exactly as in Mihatsch's Example 2.5.  Thus
\[
               \pi_*\bigl(\chi(t)t\,D\max(\varphi-t,0)\bigr)
                                  =\chi(\varphi)\varphi.
\]
The affine-cell equality extends across walls as an equality of
actual scalar functions; the right side is continuous and the
Euclidean pushforward is a delta-form of bidegree zero.

The affine projection formula now permits an arbitrary delta-form
$C$:
\[
 \pi_*\bigl(\chi(t)t\,DH\wedge\pi^*C\bigr)
       =\pi_*(\chi(t)t\,DH)\wedge C.
\]
Apply Lemma~\ref{r11p:product-integration} after moment refinement.
Only after actual moment pullback do we use the cutoff choice:
$\chi(\varphi(\ell))\varphi(\ell)=\chi(u)u=u$.
The ambient extension $\varphi$ need not be bounded, so
$\chi(\varphi)\varphi=\varphi$ is not asserted on its entire
Euclidean domain.
This argument has not pulled $C$ through the nonlinear PL graph;
only the affine projection pullback and projection formula are used.
\end{proof}

\subsection{The minimal finite trace theorem}

\begin{theorem}[Finite trace on the scalar-curvature module]
\label{r11p:minimal-trace-theorem}
Let $q:Y\to X$ be finite dominant between normal integral separated
varieties over $K$.  Every $\omega\in\mathscr M^k(Y^{\rm an})$ has a
unique actual trace on $X^{\rm an}$ satisfying
\[
                 \int_X\operatorname{Tr}^{\mathscr M}_q(\omega)
                                      \wedge\eta
                          =\int_Y\omega\wedge q^*\eta
\tag{MT9}\label{r11p:minimal-pairing}
\]
for all compact complementary actual tests.  It is local on the
ordinary topology, respects actual relations, and retains the full
native generic degree.  For a locally presented generator
$\omega=u_0Du_1\cdots Du_k$, a local formula is
\[
 \boxed{\displaystyle
 \operatorname{Tr}^{\mathscr M}_q(\omega)
 = I_{p_X}\!\left(
     \chi(t)t\,
     \operatorname{Tr}^{\rm curv}_{q\times\operatorname{id}}
       \left(D\max(u_0-t,0)\wedge\prod_{i=1}^kDu_i\right)
                         \right).}
\tag{MT10}\label{r11p:trace-formula}
\]
\end{theorem}

\begin{proof}
Shrink around an entire finite source fiber so that finitely many
generator presentations are combined as in
Lemma~\ref{r11p:pl-descent}, and $u_0$ is bounded there.  Choose
$\chi=1$ near that whole range.  The curvature trace inside
\eqref{r11p:trace-formula} exists by Lemma~\ref{r11p:pure-trace}.
Its parameter-uniform assertion supplies a product-moment presentation,
so multiplication by $\chi(t)t$ and application of $I_{p_X}$ are
legitimate.

Pair with an arbitrary compact actual test $\eta$.  First
Lemma~\ref{r11p:product-integration}, then
\eqref{r11p:pure-pairing} with the actual downstairs multiplier
$\chi(t)t\,p_X^*\eta$, and finally
Lemma~\ref{r11p:one-parameter}, give
\[
\begin{split}
 \int_X\operatorname{Tr}^{\mathscr M}_q(\omega)\eta
 &=\int_{X\times\mathbf G_m}\chi(t)t\,
     \operatorname{Tr}^{\rm curv}_{q\times\operatorname{id}}
        (DH\prod_iDu_i)\,p_X^*\eta\\
 &=\int_{Y\times\mathbf G_m}\chi(t)t\,DH\prod_iDu_i\,
                                          p_Y^*q^*\eta\\
 &=\int_Yu_0\prod_iDu_i\,q^*\eta.
\end{split}
\]
All these integrals have compact paired support.  The identity is
proved for all actual tests, not merely for model or original-smooth
tests.  Add generators.  If an expression represents zero as an
actual source form, its trace has zero pairing with every test, and
is zero by Proposition 4.11.  The same argument removes dependence
on the cutoff, source moments, quotient moments, Galois closure, and
all local choices.  It also glues the local answers in the ordinary
topology and proves uniqueness.
\end{proof}




% Source: research-sources/green-draft/supported_green.tex
\section{Supported perfect complexes and exact-support Green forms}
\label{r11p:supported-complex}

Let $K$ be complete discretely valued of characteristic zero, $R=K^\circ$,
and let $\mathcal X/R$ be integral, flat and projective, with smooth
integral projective generic fiber $X$ of dimension $n$.
The special fiber may be singular or nonreduced.  Let
\[
       C=(0\longrightarrow C_N\longrightarrow\cdots
                             \longrightarrow C_0\longrightarrow0)
\]
be a bounded complex of model vector bundles, exact on
$\mathcal U=\mathcal X\setminus\mathcal Z$, with generic open
$U=X\setminus Z$.  Assume $Z$ is pure of codimension $r$, $1\leq r\leq n$.
The same construction is linear in finite rational combinations of
such complexes.  Define
\[
 \Omega_C=\sum_i(-1)^i\operatorname{ch}_r(\overline C_i)
                              \in\mathcal B^{0,0,r}(X^{\rm an}),
 \qquad
 z_C=\operatorname{ch}^{Z}_r(C|_X)\cap[X]
                              \in\operatorname{CH}_{n-r}(Z)_{\mathbb Q}.
\tag{PC1}\label{r11p:omega-and-cycle}
\]
The first expression is an actual form by Theorem~\ref{r10c:chern}.
The second is a cycle with coefficients, because
$\operatorname{CH}_{n-r}(Z)$ is the free group on the top-dimensional
components of $Z$.  It is not merely an ambient Chow class.

\subsection{Finite multisections of flag bundles}

\begin{lemma}[A smooth finite flat multisection]
\label{r11p:multisection}
Let $F\to Y$ be a smooth projective surjection of smooth projective
integral varieties over the infinite field $K$, of pure relative
dimension $d$.  There is a smooth projective integral closed subscheme
$T_0\subset F$ whose map to $Y$ is finite flat surjective.
It may be taken as a connected component of a smooth complete intersection.
If the source is geometrically integral and $\dim Y\geq1$, $T_0$
can be chosen geometrically integral.
\end{lemma}

\begin{proof}
Choose a very ample $H$ on $F$ and use $M=H^m$ with
$m\geq\max(\dim Y,1)$.  For any positive-dimensional integral component
$A$ of any intermediate fiber complete intersection, two sections
$a_0,a_1$ of $H$ have nonconstant ratio on $A$.  Consequently
\[
                    a_0^m,a_0^{m-1}a_1,\ldots,a_1^m
\]
are linearly independent on $A$.  In the parameter space
$V=H^0(F,M)$, the sections vanishing identically on $A$ therefore
have codimension at least $m+1$.

Fix a geometric fiber, which is pure Cohen--Macaulay of dimension $d$.
Induct on $j\leq d$ to see that the bad $j$-tuples, whose common
zero locus has dimension greater than $d-j$, have codimension at least
$m+1$ in $V^j$.  Above a good prefix, a bad next section vanishes
on one of the finitely many components just considered; above a bad
prefix the inductive bound applies.  Upper semicontinuity of fiber
dimension makes the universal bad locus closed.  Adding $\dim Y$
still leaves its dimension strictly smaller than $\dim V^d$.
Properness of $Y$ makes its projection to $V^d$ closed.  Thus a
nonempty open of $d$-tuples has zero-dimensional fibers everywhere.
The fibers are nonempty, because $d$ ample divisors meet every
nonempty pure $d$-dimensional projective fiber.

Smooth Bertini supplies a second nonempty open on which the total
intersection is smooth of dimension $\dim Y$.  Geometric irreducibility
Bertini supplies the stated refinement when relevant.  Their intersection
has a $K$-rational point since $K$ is infinite.  The resulting proper
quasi-finite morphism is finite and surjective.  Its source is
Cohen--Macaulay and its target regular of the same dimension, so
miracle flatness applies.  If necessary choose a connected component;
it is integral, finite flat, and still surjective.
\end{proof}

The bound $m+1>\dim Y$, not just Bertini, is what rules out a
positive-dimensional fiber anywhere on the base.  The primary
ingredients are Stacks Tags 0D4I, 0FD6, 0G4C and 00R4.

\begin{lemma}[Relative diagonalization with the original open fixed]
\label{r11p:relative-lines}
There is a projective morphism of integral flat models
$f:\mathcal T\to\mathcal X$ with smooth projective generic fiber $T$,
such that $f:T\to X$ is generically finite of degree $d$ and finite
flat over the \emph{entire} prescribed open $U$.  There are model
line maps $s_\alpha:L_\alpha\to M_\alpha$, isomorphisms above
$\mathcal U$, and signs $\epsilon_\alpha$, satisfying
\[
                 [f^*C]
                   =\sum_\alpha\epsilon_\alpha[L_\alpha\to M_\alpha]
\tag{PC2}\label{r11p:supported-k-equality}
\]
in supported perfect $K_0$.  Their generic divisors
$D_\alpha=\operatorname{div}(s_\alpha|_T)$ are supported in $f^{-1}Z$.
\end{lemma}

\begin{proof}
On $\mathcal U$, exactness of a bounded vector-bundle complex is locally
split.  Hence $S_i^\circ=\ker d_i=\operatorname{im}d_{i+1}$ are
subbundles, with $S_0^\circ=C_0$ and $S_N^\circ=0$.  Close the graph
of this section in the product of the corresponding Grassmann bundles.
This gives a projective birational integral model
$\mathcal Y\to\mathcal X$, an isomorphism over $\mathcal U$, with
tautological subbundles $S_i\subset C_i$ and quotient bundles $Q_i$.
Every asserted vanishing of a bundle map that holds on the dense open
holds on this integral model.  The differential therefore factors as
\[
                            a_i:Q_i\longrightarrow S_{i-1},
\]
an equal-rank map that is an isomorphism over $\mathcal U$.

There is an actual filtration by subcomplexes: its $p$th term has
$C_i$ for $i>p$, $S_p$ in degree $p$, and zero below $p$.
The successive two-term quotients are
$Q_{p+1}\to S_p$ in degrees $p+1,p$.  They are exact on the given
open, and yield
\[
                     [C]=\sum_i(-1)^{i-1}[Q_i\to S_{i-1}]
\]
in supported $K_0$, rather than only in ordinary $K_0$.

Resolve the generic $\mathcal Y_K$ in characteristic zero, without
changing its smooth open $U$, and extend the projective resolution by
graph closure over $R$.  This uses resolution over the field $K$ only,
not mixed-characteristic resolution of a model.  Call the resulting
smooth projective generic fiber $Y$.

Take the product of the full source flag bundles of \emph{all} $Q_i$
over $Y$.  Before extending the transported target flags, choose a
smooth finite flat multisection $T_0\to Y$ by
Lemma~\ref{r11p:multisection}.  Take its flat model closure.  Over
$\mathcal U$, the isomorphisms $a_i$ transport every source flag to a
target flag.  Close the graph of these target flags, and resolve the
resulting generic fiber without changing its smooth inverse image of
$U$.  Its final generic fiber $T$ is therefore still exactly $T_0$
over $U$.  In particular $T\to X$ is finite flat over all of $U$;
the successive modifications have introduced no new excluded image.

The maps $a_i$ preserve the two full flags.  Their associated graded
maps are model line maps $L_\alpha\to M_\alpha$, isomorphisms on the
original model open.  Termwise filtration of the two-term complexes
gives \eqref{r11p:supported-k-equality}.  A nonzero map of line bundles
on the integral model is injective and defines an effective Cartier
divisor.  Its generic divisor can only lie above $Z$.  Empty generic
divisors are allowed, for instance for a purely vertical line map.
\end{proof}

\subsection{Localized characters and cycle coefficients}

For degree-one variables $l,m$ put
\[
 B_{r-1}(l,m)=\frac1{r!}\sum_{a=0}^{r-1}m^{r-1-a}l^a,
 \qquad
                  (m-l)B_{r-1}(l,m)=\frac{m^r-l^r}{r!}.
\tag{PC3}\label{r11p:B-polynomial}
\]
For a nonzero line map $s:L\to M$ on $T$, with divisor $D$, the
all-degree Cartier localized character formula is
\[
 \operatorname{ch}^{D}_r([L\to M])\cap[T]
         =B_{r-1}(c_1L,c_1M)\cap[D].
\tag{PC4}\label{r11p:cartier-localized}
\]
Here $L$ has degree one and $M$ degree zero.  To check the full
polynomial, the complex resolves $M|_D$ and the normal bundle of
$D\hookrightarrow T$ is $(M\otimes L^{-1})|_D$.  Localized
Riemann--Roch for this regular immersion gives
\[
 e^m\,\frac{1-e^{-(m-l)}}{m-l}\cap[D]
                         =\frac{e^m-e^l}{m-l}\cap[D].
\]
Its homogeneous degree $r-1$ is exactly \eqref{r11p:cartier-localized}.
This is Fulton, \emph{Intersection Theory}, Corollary 18.1.2,
specialized to a Cartier immersion; the proof uses deformation to the
normal bundle.  A leading-term statement alone, such as Stacks 0FEH,
would not justify this all-degree polynomial.

Additivity on supported perfect $K_0$, compatibility with arbitrary
base change, and the proper bivariant projection formula now give
\[
 \sum_\alpha\epsilon_\alpha
    (f|_{D_\alpha})_*
       \left(B_{r-1}(c_1L_\alpha,c_1M_\alpha)\cap[D_\alpha]\right)
                              =d\,z_C
       \quad\hbox{in }\operatorname{CH}_{n-r}(Z)_{\mathbb Q}.
\tag{PC5}\label{r11p:source-equality}
\]
The underlying equality is
$f_*(\operatorname{ch}^{f^{-1}Z}(f^*C)\cap[T])
=\operatorname{ch}^Z(C)\cap f_*[T]$, and $f_*[T]=d[X]$.
These are supported identities; one has not first forgotten the
support and then attempted to recover coefficients.  Nilpotent
thickenings can be used to make every indicated factorization through
$Z$ scheme-theoretic.  Chow groups are unchanged by those thickenings.
Since the target Chow group is top-dimensional on $Z$,
\eqref{r11p:source-equality} is equality of actual cycle coefficients.
See Stacks 0FB3 (base change), 0FBE (additivity), and 0FET (supported
perfect $K_0$).

\subsection{Chern degrees and actual top-degree tests}

\begin{lemma}[Chern degree against an arbitrary actual top test]
\label{r11p:proper-chern-degree}
Let $h:B\to A$ be proper between integral projective $K$-varieties,
$\dim A=m$, $\dim B=m+t$.  Let $P$ be a homogeneous degree-$t$
polynomial in first Chern forms of algebraic line bundles on $B$
with PWS metrics.  For every compact actual top form $\eta$ on $A$,
\[
              \int_BP\wedge h^*\eta
                =\int_{h_*(P_{\rm alg}\cap[B])}\eta.
\tag{PC6}\label{r11p:proper-top-pairing}
\]
The statement is linear in cycles and includes their generic lengths.
\end{lemma}

\begin{proof}
Write each algebraic line bundle as a quotient of two very ample ones.
Choose PWS metrics on those factors so that their quotient is the
given metric.  Expand $P$ by multilinearity.  For each resulting
monomial, choose a section of each very ample factor successively,
meeting the finitely many current integral cycle components properly.
The infinite field $K$ supplies such choices.  Work componentwise
when an intermediate divisor is reducible or nonreduced.

At every step apply scalar strong Poincare--Lelong, Mihatsch
Corollary 5.6, with the remaining closed curvature factors and the
pullback of $\eta$ as test.  The latter is separately closed because
$\eta$ is top-degree on $A$.  The Green correction is therefore zero.
One is left with the pullback of $\eta$ integrated on a cycle of
dimension $m$.

A component mapping to a lower-dimensional subset of $A$ contributes
zero, because the support of the actual top test avoids that subset.
For a dominating component, remove a proper closed subset of $A$
so that the component maps finite flat of its native generic degree.
The top test again avoids the removed subset.  The native finite
degree formula evaluates the integral with exactly that degree.
This is precisely algebraic proper pushforward of the terminal
intersection cycle, proving \eqref{r11p:proper-top-pairing}.
\end{proof}

This scalar equality applies to
each $D_\alpha\to Z_j$ in \eqref{r11p:source-equality}, with
$m=n-r$ and $t=r-1$.  Components whose image has smaller dimension
contribute zero by the same support argument.

\subsection{The exact-open Green construction}

\begin{theorem}[A strong Green pair for a supported model complex]
\label{r11p:complex-green}
Under the assumptions of \eqref{r11p:omega-and-cycle}, there is an
actual form
\[
                       G_C\in\mathcal B^{r-1,r-1}(U^{\rm an})
\]
such that, for every compact actual
$\eta\in\mathcal B_c^{n-r,n-r}(X^{\rm an})$,
\[
                 \int_UG_C\wedge D\eta+\int_{z_C}\eta
                                     =\int_X\Omega_C\wedge\eta.
\tag{PC7}\label{r11p:complex-green-equation}
\]
The domain is the prescribed $U=X\setminus Z$, with no auxiliary
closed subset removed.  If a tube is disjoint from the full model
support $\mathcal Z$, the construction can be chosen zero on that tube.
\end{theorem}

\begin{proof}
Use Lemma~\ref{r11p:relative-lines} and give the line bundles their
model metrics.  Write
\[
 l_\alpha=c_1(\overline L_\alpha),\qquad
 m_\alpha=c_1(\overline M_\alpha),\qquad
 g_\alpha=-\log\|s_\alpha\|.
\]
On $f^{-1}U$ each $g_\alpha$ is a PL scalar.  Put
\[
       G_T=\sum_\alpha\epsilon_\alpha
                    g_\alpha B_{r-1}(l_\alpha,m_\alpha),
       \qquad
       G_C=\frac1d\operatorname{Tr}^{\mathscr M}_{f|_{f^{-1}U}}(G_T).
\tag{PC8}\label{r11p:complex-green-formula}
\]
Every summand lies in the scalar-curvature module of
Theorem~\ref{r11p:minimal-trace-theorem}.  Source and target of the
finite restriction are smooth, hence normal; all its assumptions hold.
The resulting form is actual on the full ordinary open $U$.

Strict model Chern additivity and functoriality from Theorem~\ref{r10c:chern} convert
\eqref{r11p:supported-k-equality}, after forgetting support, into
the actual identity on $T$
\[
 \sum_\alpha\epsilon_\alpha
             (m_\alpha-l_\alpha)B_{r-1}(l_\alpha,m_\alpha)
                                      =f^*\Omega_C.
\tag{PC9}\label{r11p:actual-curvature-equality}
\]


Apply scalar strong Poincare--Lelong to each $s_\alpha$ with test
$B_{r-1}(l_\alpha,m_\alpha)f^*\eta$.  These tests are compact because
$f$ is proper.  Curvature factors are separately closed.  Summation
and \eqref{r11p:actual-curvature-equality} yield
\[
 \int_{T\setminus f^{-1}Z}G_Tf^*D\eta
 +\sum_\alpha\epsilon_\alpha
       \int_{[D_\alpha]}B_{r-1}(l_\alpha,m_\alpha)f^*\eta
                         =\int_Tf^*(\Omega_C\eta).
\tag{PC10}\label{r11p:upstairs-complex-PL}
\]
For the first term, extending each off-divisor integral to the
displayed common open changes nothing: $D\eta$ has bidegree
$(n-r+1,n-r+1)$, so its closed support avoids $Z$ by
Mihatsch Corollary 4.9.  Properness puts its pulled-back support
inside the finite locus $f^{-1}U$.  The minimal trace pairing identifies
this first term with $d\int_UG_CD\eta$.

The form $\Omega_C\eta$ is top-degree on $X$, so its support also
avoids $Z$.  On that support $f$ is finite flat of degree $d$;
the last term is $d\int_X\Omega_C\eta$ by the native degree formula.
The middle term equals $d\int_{z_C}\eta$ by
Lemma~\ref{r11p:proper-chern-degree} and
\eqref{r11p:source-equality}.  Divide by $d$.

All tests defining the off-$Z$ integration functional of $G_C$
have bidegree $(n-r+1,n-r+1)$ and therefore support away from $Z$.
Thus its canonical extension and the all-actual meaning of
\eqref{r11p:complex-green-equation} require no integrability
assumption near $Z$ beyond those already encoded by support avoidance.

Finally, above the model open $\mathcal U$ each $s_\alpha$ is a model
isomorphism of line bundles, so its model norm is one and $g_\alpha=0$.
The local trace then vanishes over every tube avoiding $\mathcal Z$.
This concerns the full model support, including possible vertical
cohomology, not just the closure of its generic support.
\end{proof}

\paragraph{Support of the prescribed representative.}
The source of a supported complex can have more components than its
associated cycle after coefficient cancellation.  Choosing one complex
for each nonzero integral component of a prescribed cycle gives the
form on its exact complement, as follows.

\subsection{Arbitrary prescribed cycles and the good-reduction sector}

\begin{corollary}[Every prescribed cycle has a strong model-curvature pair]
\label{r11p:any-cycle-pair}
Let $X/K$ be smooth geometrically integral projective and let
$Z=\sum_j a_j[Z_j]$ be a rational codimension-$r$ cycle, written with
distinct integral components and nonzero coefficients.  After choosing
projective model data, there are an actual model Chern form $\Omega_Z$
and an actual
\[
                  G_Z\in\mathcal B^{r-1,r-1}
                           (X^{\rm an}\setminus|Z|^{\rm an})
\]
satisfying $D[G_Z]_\delta+\delta_Z=[\Omega_Z]_\delta$ on all compact
actual tests.  No generic cycle representative is moved.
\end{corollary}

\begin{proof}
On the regular projective generic fiber, $\mathcal O_{Z_j}$ has a
finite globally locally free resolution.  Its degree-$r$ supported
character is $[Z_j]$, including the leading coefficient; this is
Stacks 0FEW, or the leading-term part of localized Riemann--Roch.
It is not necessary for $Z_j$ to be a regular immersion.

These finitely many generic bundle complexes extend to a common
projective model after a vertical modification.  For completeness,
extend each bundle coherently, twist to obtain finitely many generators,
and close the graph of its generic quotient in a Grassmann bundle.
Its tautological quotient is a model vector bundle with the original
generic fiber.  Use the product of the finitely many Grassmann bundles.
Every generic differential extends after multiplying by a common
power of a uniformizer; the resulting successive products are zero
because they vanish on the dense generic fiber.  The scaled generic
complex is chain-isomorphic to the original one by degreewise scalar
rescaling.  Its full model cohomology support has generic support
exactly $Z_j$.

Apply Theorem~\ref{r11p:complex-green} to each complex and sum with
coefficients $a_j$.  The union of the removed sets is exactly
$|Z|$, because the prescribed cycle was first written with distinct
nonzero components.  The zero cycle is handled by the zero form.
\end{proof}

\begin{corollary}[Strong G for arbitrary cycles on a smooth good model]
\label{r11p:good-reduction-G}
Suppose $X$ admits a smooth projective $R$-model $\mathcal X$.
For every rational codimension-$r$ cycle $Z$ on $X$ whose geometric
$\ell$-adic class is zero, $\ell\ne\operatorname{char}(k)$, the form
in Corollary~\ref{r11p:any-cycle-pair} can be chosen with
$\Omega_Z=0$:
\[
                  \int_{X\setminus|Z|}G_ZD\eta=-\int_Z\eta
       \qquad
       \bigl(\eta\in\mathcal B_c^{n-r,n-r}(X^{\rm an})\bigr).
\tag{PC11}\label{r11p:good-reduction-zero}
\]
This includes arbitrary singular or non-complete-intersection cycle
components and retains the exact support and original field.
\end{corollary}

\begin{proof}
For each prescribed component take its scheme-theoretic flat closure
$\mathcal Z_j\subset\mathcal X$.  The smooth model over a DVR is regular
of finite dimension and has an ample line bundle.  Hence
$\mathcal O_{\mathcal Z_j}$ has a finite globally locally free
resolution on this fixed model.  One explicit justification is to
surject finite sums of negative ample twists onto a coherent sheaf
and its successive coherent kernels.  After at most $\dim\mathcal X$
steps the remaining kernel is locally free, since every regular
local ring has global dimension its Krull dimension.  Thus the
global assertion is not inferred from perfectness without using the
resolution property.  No modification of the good model
is needed.  Use those model complexes in the preceding construction.

Their alternating degree-$r$ Chern character on the generic fiber
is $[Z]$ in $\operatorname{CH}^r(X)_{\mathbb Q}$.  Smooth proper
base change and compatibility of etale Chern classes identify its
geometric special-fiber Chern class with the zero generic class.
Its degree against every integral $r$-cycle on the smooth special
fiber is therefore zero by the etale trace pairing.

Theorem~\ref{r10c:numerical-zero} now gives
$\Omega_Z=0$ as an actual form, not just as a current or a
cohomology class.  Substitute into the strong equation already
proved against every actual test.
\end{proof}

\begin{corollary}[Finite good covers and potential good reduction]
\label{r11p:potential-good-G}
Let $X/K$ be smooth integral projective.  Suppose there are a finite
valued extension $K'/K$, a smooth integral projective $K'$-variety
$Y$ with a smooth projective $R'=K'^\circ$-model, and a finite
dominant $K$-morphism $h:Y\to X$.  Then every rational geometrically
homologically trivial codimension-$r$ cycle $Z$ on $X$ admits a
zero-curvature strong Green form on exactly
$X^{\rm an}\setminus|Z|^{\rm an}$, over the original field $K$,
against every compact $K$-actual complementary test.

In particular this holds when $X$ has potential good reduction:
take $Y=X_{K'}$ and its finite projection to $X$, working componentwise
if the base change is disconnected.
It also holds for a finite dominant smooth projective good-reduction
cover over $K$ itself.
\end{corollary}

\begin{proof}
Regard $Y$ as a $K$-variety.  It is smooth and integral, though it
need not be geometrically integral over $K$.  This is enough for
the preceding supported-complex theorem and the minimal trace:
the multisection proof permits choosing a smooth connected component
finite flat over an integral base.  Since $Y$ is Cohen--Macaulay
and $X$ is regular of the same dimension, the finite dominant $h$
is flat by miracle flatness.  Write $a=\deg_K h$ and set $Z'=h^*Z$,
with the scheme-theoretic generic lengths in flat pullback.
Its support is exactly $h^{-1}|Z|$, its geometric cohomology class
is zero, and
\[
                                  h_*Z'=aZ.
\tag{PC12}\label{r11p:finite-cover-cycle-degree}
\]

Resolve the closures of the distinct components of $Z'$ on the
fixed smooth $R'$-model $\mathcal Y$, with their rational
coefficients.  View $\mathcal Y$ also as a projective flat
$R$-model of the $K$-variety $Y$.  If $e$ is the ramification index,
its special fiber over $R$ is the $e$th nilpotent thickening of its
good special fiber over $R'$: locally a uniformizer of $R$ is a
unit times the $e$th power of one of $R'$.

Every integral $r$-cycle on this thickened special fiber lies on
its reduction.  The restriction of the alternating operational
degree-$r$ Chern class there is geometrically homologically zero
by smooth proper base change over $R'$.  Thus it has degree zero
against every such cycle.  Computing degrees over the smaller
residue field only multiplies them by $[k':k]$, so they remain zero.
Apply Theorem~\ref{r10c:numerical-zero} directly to the $R$-model
$\mathcal Y$.  It gives
\[
              \Omega_{Z'}=0
                       \quad\hbox{as a $K$-actual form on $Y^{\rm an}$}.
\]
Native special-fiber multiplicities, including $e$, are retained;
none has been replaced by one.  This step does not infer the
$K$-actual statement from equality on a smaller collection of
$K'$-tests.

Perform the relative line construction over this $R$-model
simultaneously for the finitely many chosen resolutions.  Close the
joint kernel-Grassmann graph over their common exact open, then take
the product of all source flag bundles before choosing one smooth
finite multisection.  Transport and close all target flags on that
same multisection.  This gives one common map carrying the rationally
weighted sum of all the line complexes; it does not use a fiber
product of separately chosen branched finite multisections.
Keep its line sum $G_T$ before applying any trace.  Let $b:T\to Y$
be the resulting map of generic degree $d$, finite flat over
$Y\setminus h^{-1}|Z|$.  The composite
$f=h\circ b:T\to X$ is finite flat of degree $ad$ above the whole
$U=X\setminus|Z|$.  Define directly over $K$
\[
                   G_Z=\frac1{ad}
                       \operatorname{Tr}^{\mathscr M}_{f|_{f^{-1}U}}
                                      (G_T|_{f^{-1}U}).
\tag{PC13}\label{r11p:potential-good-formula}
\]
This is a permitted trace: $G_T$ is the original finite sum of
PL scalars times model curvature polynomials.  We have not assumed
that an already constructed $G_{Z'}$ lies in the trace module.

Repeat the scalar Poincare--Lelong proof of
Theorem~\ref{r11p:complex-green} with an arbitrary compact actual
test $\eta$ on $X$.  The source curvature is
$b^*\Omega_{Z'}=0$.  Its residue pushed to $Y$ is $dZ'$, hence
pushed to $X$ is $daZ$ by
\eqref{r11p:finite-cover-cycle-degree}.  The support of $D\eta$
avoids $|Z|$, so the first term is evaluated entirely on the
finite open for the composite $f$.  Its minimal-trace pairing
is $ad\int_U G_ZD\eta$.  Division by $ad$ proves
$\int_U G_ZD\eta=-\int_Z\eta$ on all the stated tests and on the
exact final complement.

For potential good reduction with geometrically integral $X$, the
native degree of the projection $X_{K'}\to X$ is $[K':K]$.
If that base change is disconnected, a smooth integral connected
component and its good-model component may instead be used, with its
actual native degree over $X$.  The possible ramification factor is
already part of the degree and of the model calibrations; no
additional normalization is inserted.
\end{proof}

\subsection{Global primitives for generic Chow-zero model characters}

\begin{corollary}[Generically acyclic complexes]
\label{r11p:acyclic}
If a bounded model complex is exact on all of $X$, then
$\sum_i(-1)^i\operatorname{ch}_r(\overline C_i)$ has a global actual
$D$-primitive on $X^{\rm an}$.
\end{corollary}

\begin{proof}
The same relative construction is made with generic open $U=X$;
its final generic multisection is finite flat over all of $X$ and all
generic line divisors are empty.  Formula \eqref{r11p:complex-green-formula}
is global, and \eqref{r11p:complex-green-equation} has zero residue.
Actual current injectivity identifies its $D$ with the displayed
actual curvature.
\end{proof}

\begin{proposition}[A generic $K_0$ relation has an actual primitive]
\label{r11p:generic-k-primitive}
Let $\xi\in K_0(\mathcal X)_{\mathbb Q}$.  If $\xi|_X=0$ in
$K_0(X)_{\mathbb Q}$, then every positive-degree actual
$\operatorname{ch}_r(\xi)$ has a global actual $D$-primitive.
\end{proposition}

\begin{proof}
A relation in $K_0(X)$ is a finite combination of short exact sequence
relations and differences of isomorphic vector bundles.  Extend the
finitely many bundles to a common vertical model refinement, which
does not change any original model metric on the generic fiber.
For a generic exact sequence $0\to A\to B\to C\to0$, close the graph
of the subbundle $A\subset B$ in a Grassmann bundle of the already
chosen model extension of $B$.  This produces a model exact sequence
\[
                         0\to S\to B\to Q\to0
\]
with $S|_X\simeq A$ and $Q|_X\simeq C$.  Strict model Whitney gives
\[
      \operatorname{ch}_r(B-A-C)
       =\operatorname{ch}_r(S-A)+\operatorname{ch}_r(Q-C).
\]
After clearing uniformizer powers, the two generic isomorphisms are
model maps giving two-term generically acyclic complexes.  Their
global primitives exist by Corollary~\ref{r11p:acyclic}.
The same handles the isomorphic-bundle differences in the chosen
finite presentation of the $K_0$ relation.  Add the primitives and
divide by the cleared integer denominator.
\end{proof}

\begin{proposition}[Only the generic degree-$r$ Chow class matters]
\label{r11p:generic-chow-primitive}
If $\xi\in K_0(\mathcal X)_{\mathbb Q}$ satisfies
$\operatorname{ch}_r(\xi|_X)=0$ in
$\operatorname{CH}^r(X)_{\mathbb Q}$, then its actual model
$\operatorname{ch}_r(\xi)$ has a global actual $D$-primitive.
\end{proposition}

\begin{proof}
Let $\psi^2$ be the Adams operation on model vector-bundle $K_0$ and set
\[
 e_r(T)=\prod_{\substack{0\leq j\leq n\\j\ne r}}
                  \frac{T-2^j}{2^r-2^j},
 \qquad
                            \xi_r=e_r(\psi^2)\xi.
\]
On the smooth projective generic fiber the rational Chern character
$K_0(X)_{\mathbb Q}\simeq\bigoplus_{j=0}^n
\operatorname{CH}^j(X)_{\mathbb Q}$ is an isomorphism, and
$\operatorname{ch}_j(\psi^2\zeta)=2^j\operatorname{ch}_j(\zeta)$.
All generic Chern components of $\xi_r$ are zero, so $\xi_r|_X=0$.

The same Adams identity is strict for actual model Chern forms:
it is a universal symmetric-root identity, hence follows from
Theorem~\ref{r10c:chern}.
Therefore $\operatorname{ch}_r(\xi_r)=\operatorname{ch}_r(\xi)$
as actual forms.  Apply Proposition~\ref{r11p:generic-k-primitive}
to $\xi_r$.
\end{proof}

\begin{corollary}[Chow-zero cycles without a reduction hypothesis]
\label{r11p:chow-zero-G}
For every prescribed rational codimension-$r$ cycle $Z$ with
$[Z]=0$ in $\operatorname{CH}^r(X)_{\mathbb Q}$, strong G holds
on exactly $X^{\rm an}\setminus|Z|^{\rm an}$, over $K$, without
a good-reduction assumption.
\end{corollary}

\begin{proof}
Construct the exact-support pair $(G_Z,\Omega_Z)$ above.
Its model curvature is $\operatorname{ch}_r(\xi)$ for a model
virtual class with generic degree-$r$ character $[Z]=0$.
Choose a global actual $\Psi$ with $D\Psi=\Omega_Z$ by the preceding
proposition.  Then $G_Z-\Psi|_{X\setminus|Z|}$ has zero curvature
and the same source cycle, with no new poles.
\end{proof}




% Source: research-sources/green-draft/algebraic_equivalence.tex
\section{Algebraic equivalence and compact parameters}
\label{r12p:section}

Throughout this section $K/\mathbf Q_p$ is finite and $X/K$ is smooth,
projective and geometrically integral.  Write $D=d'd''$ and let
$\mathcal B$ be the actual delta-form sheaf.

Recall the algebraic quotient
\[
 H_{\mathrm{BC,act}}^{r,r}(X)=
 \frac{\{\omega\in\mathcal B^{r,r}(X^{\rm an}):
                    d'\omega=d''\omega=0\}}
      {D\mathcal B^{r-1,r-1}(X^{\rm an})}.
\]
The preceding theorems define the model-independent class
\[
 \mathfrak o_X^r:\operatorname{CH}^r(X)_{\mathbb Q}
       \longrightarrow H_{\mathrm{BC,act}}^{r,r}(X),
 \qquad
 \mathfrak o_X^r(a)=[\operatorname{ch}_r(\bar\xi)],
 \quad \operatorname{ch}_r(\xi|_X)=a.
\tag{R12.P1}\label{r12p:obstruction}
\]


\subsection{Actual forms are locally constant in a rigid parameter}

\begin{lemma}[Uniform rigid-fiber constancy]
\label{r12p:rigid-fiber}
Let $S/K$ be smooth and separated, let $a_0\in S(K)$, and let
$\omega\in\mathcal B((X\times S)^{\rm an})$.  Denote by
$s_a:X^{\rm an}\to(X\times S)^{\rm an}$ the section with parameter
$a\in S(K)$.  There is an ordinary analytic open neighborhood $V$ of
$a_0$ in $S^{\rm an}$ such that
\[
                    s_a^*\omega=s_{a_0}^*\omega
       \quad\text{for every }a\in V\cap S(K),
\tag{R12.P2}\label{r12p:literal-constancy}
\]
as actual forms on the whole $X^{\rm an}$.
\end{lemma}

\begin{proof}
Put $p:(X\times S)^{\rm an}\to X^{\rm an}$ and
$q:(X\times S)^{\rm an}\to S^{\rm an}$.  The analytic section
$s=s_{a_0}$ is closed and $ps=\operatorname{id}$.  We first check
the retracted germ argument directly.

Near a point of $s(X^{\rm an})$, write
$\omega=t_f^\star\gamma$ for a finite tuple $f_i$ of analytic units
and a Euclidean delta-form $\gamma$.  Shrink to where both $f_i$ and
$f_i^0=p^*s^*f_i$ are defined.  Each $f_i-f_i^0$ vanishes along the
section.  The inequalities
\[
                         |f_i-f_i^0|<|f_i^0|
\]
therefore define an ordinary open neighborhood of that part of the
section.  On it $|f_i|=|f_i^0|$.  On every common refining skeleton
the corresponding valuation tuples are equal as linear functions,
so their intrinsic tropical pullbacks of $\gamma$ agree.  Mihatsch's
Definition 4.2 and Lemma 4.4 identify the actual forms, and Theorem 4.5
gives the compatibility with analytic pullback \cite{M}.  Taking the
union of the resulting neighborhoods gives an ordinary open $W$
containing the entire section, on which
\[
                         \omega=p^*s^*\omega.
\]

The projection $q$ is proper because $X$ is proper.  In particular it
is closed on the ordinary topology.  Its closed set
$q((X\times S)^{\rm an}\setminus W)$ does not contain $a_0$:
the analytic fiber over a $K$-rational point is exactly the image of
the corresponding section.  Its complement $V$ is an ordinary open
with $q^{-1}(V)\subset W$.  Pull the displayed equality back by
$s_a$ for $a\in V\cap S(K)$ to obtain
\eqref{r12p:literal-constancy}.

This argument uses properness of the analytic projection, not an
identification of the underlying analytic product with a Cartesian
product of point sets.  All pullbacks are of actual forms by genuine
analytic morphisms, including the closed section.
\end{proof}

\begin{lemma}[A multiple enters a prescribed identity neighborhood]
\label{r12p:compact-multiple}
Let $A/K$ be an abelian variety.  For every $a\in A(K)$ and every
open neighborhood $U$ of $0$ in $A(K)$, there is an integer $m>0$
such that $ma\in U$.
\end{lemma}

\begin{proof}
The group $A(K)$ is compact: a projective embedding realizes it as
a closed subset of compact $\mathbf P^N(K)$.  Choose a neighborhood
$V$ of $0$ with $V-V\subset U$.  Finitely many translates of $V$
cover $A(K)$.  Two of sufficiently many points $0,a,2a,\ldots,Na$
belong to one translate; their difference is $ma\in U$ for some
$m>0$.  This remains valid when $a$ is torsion.
\end{proof}

\subsection{Additive parameter families and global primitives}

\begin{definition}[Additive cycle family]
For an abelian variety $A/K$, a class
$\Theta\in\operatorname{CH}^r(X\times A)_{\mathbb Q}$ is called
additive on rational fibers if its refined Gysin fibers satisfy
$\Theta_0=0$ and
\[
       \Theta_{a+b}=\Theta_a+\Theta_b
       \quad\text{in }\operatorname{CH}^r(X)_{\mathbb Q}
       \qquad(a,b\in A(K)).
\tag{R12.P3}\label{r12p:additive}
\]
The requirement is about Chow classes.  An arbitrary family on an
abelian parameter variety is not presumed to satisfy it.
\end{definition}

\begin{theorem}[Compact additive parameters]
\label{r12p:compact-additive}
Let $\Theta$ be as in \eqref{r12p:additive}.  For every $a\in A(K)$,
\[
                          \mathfrak o_X^r(\Theta_a)=0.
\tag{R12.P4}\label{r12p:additive-zero}
\]
More precisely, any actual degree-$r$ model character representing
$\Theta_a$ has a global actual $D$-primitive that is a finite rational
linear combination of the restricted finite-flat traces of global
original PL-scalar/model-curvature line sums from Theorem~\ref{r11p:minimal-trace-theorem}.
\end{theorem}

\begin{proof}
Choose a virtual bundle on $X\times A$ whose degree-$r$ generic
Chern character is $\Theta$, and extend its finitely many bundles
to a common projective flat model.  This uses the rational Chern
character isomorphism on the smooth projective generic fiber and
the model-extension argument of Corollary~\ref{r11p:any-cycle-pair}.
Let $\Omega$ be its actual degree-$r$ model character and write
$\Omega_b=s_b^*\Omega$.  Each $\Omega_b$ is a model character on a
graph model of $X$.  For any finite list of parameters, all of these
model bundles can be compared on one projective flat refinement.

Lemma~\ref{r12p:rigid-fiber} gives a neighborhood $U$ of $0$ in
$A(K)$ on which $\Omega_b=\Omega_0$ literally.  Fix $a$, and choose
$m>0$ with $ma\in U$ by Lemma~\ref{r12p:compact-multiple}.  On a
common model the actual character
\[
 \Delta=\Omega_{ma}-\Omega_0-m(\Omega_a-\Omega_0)
                 =-m(\Omega_a-\Omega_0)
\tag{R12.P5}\label{r12p:delta}
\]
has generic degree-$r$ Chow class
$\Theta_{ma}-m\Theta_a=0$.  Proposition~\ref{r11p:generic-chow-primitive} supplies a global actual $H$ with $DH=\Delta$.
Consequently
\[
                 \beta_a=-\frac{1}{m}H,
                 \qquad D\beta_a=\Omega_a-\Omega_0.
\tag{R12.P6}\label{r12p:beta}
\]

If $\Omega_a'$ is any other actual degree-$r$ model character with
generic class $\Theta_a$, then
$\Omega_a'-(\Omega_a-\Omega_0)$ has generic class zero.  Give this
difference a global primitive by the same proposition and add it to
$\beta_a$.  This proves \eqref{r12p:additive-zero} and independence
of the chosen model metric at the level of actual primitives.

For the asserted form of the primitive, the generic-Chow-zero proof
first applies an Adams projector, then expresses a generic $K_0$
relation by finitely many generically acyclic model complexes.
Their prescribed generic open is the whole $X$.  Each primitive is
therefore of the form
\[
       \sum_j q_jd_j^{-1}\operatorname{Tr}^{\mathscr M}_{b_j}(G_j),
       \qquad q_j\in\mathbb Q,
\tag{R12.P7}\label{r12p:certificate}
\]
where $b_j:T_j\to X$ is finite flat on all of $X$, $d_j$ is its
native degree, and $G_j$ is a global original line sum on $T_j$.
Taking the finite rational linear combinations in
\eqref{r12p:beta} and in the final correction preserves this
description.  No trace of an arbitrary already constructed actual
primitive is used.
\end{proof}

\begin{remark}[A shorter quotient argument]
By Lemma~\ref{r12p:rigid-fiber}, the map
$a\mapsto\mathfrak o_X^r(\Theta_a)$ is locally constant as a map
from the compact space $A(K)$ to an abstract set.  It therefore has
finite image.  Additivity makes that image a finite subgroup of a
real vector space, hence zero.  This proves the vanishing without
putting any topology on $H_{\mathrm{BC,act}}^{r,r}(X)$.
The proof above additionally retains the composable representation
\eqref{r12p:certificate}.
\end{remark}

\subsection{The Poincare correspondence and descent over the original field}

\begin{lemma}[An additive family attached to a pointed curve]
\label{r12p:poincare}
Let $F/K$ be a finite extension, let $C/F$ be a smooth projective
geometrically integral curve with $c_0\in C(F)$, and let
$\Gamma\in\operatorname{CH}^r(X_F\times C)_{\mathbb Q}$.
Put $J=\operatorname{Pic}^0_{C/F}$.  There is a class
$\Theta\in\operatorname{CH}^r(X_F\times J)_{\mathbb Q}$ such that
\[
 \Theta_t=\Gamma_*c_1(\mathcal P_t),\qquad
 \Theta_0=0,\qquad \Theta_{t+t'}=\Theta_t+\Theta_{t'}
 \quad(t,t'\in J(F)),
\tag{R12.P8}\label{r12p:poincare-fiber}
\]
where $\mathcal P$ is the Poincare bundle normalized along
$c_0\times J$ and $C\times\{0\}$.
\end{lemma}

\begin{proof}
The pointed Picard functor has this normalized universal line bundle;
see Milne's Theorem 1.2 and Lemma 1.3 \cite{MilneJV}.  On
$X_F\times C\times J$ define
\[
 \Theta=(p_{XJ})_*\bigl(p_{XC}^*\Gamma\cdot
                              p_{CJ}^*c_1(\mathcal P)\bigr).
\tag{R12.P9}\label{r12p:poincare-family}
\]
Here $p_{XJ}$ has relative dimension one, so the codimension is $r$.
Every pushforward in this formula is the usual algebraic proper
pushforward on Chow groups.  Refined Gysin compatibility with that
pushforward and with intersection gives the fiber formula.
The rigidified Picard law identifies
$\mathcal P_{t+t'}$ with $\mathcal P_t\otimes\mathcal P_{t'}$,
and $\mathcal P_0$ is trivial.  First Chern classes and correspondence
action are additive, proving all of \eqref{r12p:poincare-fiber}.
\end{proof}

\begin{lemma}[Algebraic norm of an additive family]
\label{r12p:norm-family}
Let $L/K$ be finite Galois, let $J/L$ be an abelian variety, and let
$\Theta\in\operatorname{CH}^r(X_L\times J)_{\mathbb Q}$ be
additive on $J(L)$.  Put $A=\operatorname{Res}_{L/K}J$.
There is a class
$\Theta^N\in\operatorname{CH}^r(X\times A)_{\mathbb Q}$,
additive on $A(K)$, whose fiber at the point $a$ corresponding to
$t\in J(L)=A(K)$ is
\[
                    (\Theta^N)_a=h_*\Theta_t,
                    \qquad h:X_L\longrightarrow X.
\tag{R12.P10}\label{r12p:norm-fiber}
\]
\end{lemma}

\begin{proof}
For completeness, $A$ can be constructed by descending the product
$\prod_{\sigma\in\operatorname{Gal}(L/K)}{}^\sigma J$ with its
canonical semilinear permutation datum.  Finite Galois descent for
projective varieties gives a smooth projective group, geometrically
connected, hence an abelian variety over $K$; see
\cite[Proposition 1.8]{MilneJV}.  Its $K$-points are the conjugate
tuples determined by a point of $J(L)$.  Let $e:A_L\to J$ be the
evaluation morphism, which is a group morphism.  With
$h_A:(X\times A)_L\to X\times A$, set
\[
 \Theta^N=(h_A)_*(\operatorname{id}_{X_L}\times e)^*\Theta.
\tag{R12.P11}\label{r12p:norm-family-formula}
\]
The pullback here is an algebraic Chow pullback between smooth
schemes, and $h_A$ is finite flat.  The proper-pushforward/refined-
Gysin base-change formula gives \eqref{r12p:norm-fiber}.
Since $e$ is additive and $h_*$ is a homomorphism, these fibers are
additive; their zero fiber is zero.  This proof only uses algebraic
norms.  It neither pulls actual forms back to $L$ nor descends them
through an analytic finite trace.
\end{proof}

\begin{theorem}[Geometrically algebraically trivial curvature]
\label{r12p:geometric-algebraic}
Let $1\leq r\leq n$ and let
$z\in\operatorname{CH}^r(X)_{\mathbb Q}$ become algebraically
trivial over $\overline K$.  Then
\[
                              \mathfrak o_X^r(z)=0.
\tag{R12.P12}\label{r12p:alg-zero}
\]
Every actual degree-$r$ model character representing $z$ has a
global actual $D$-primitive of the form
\eqref{r12p:certificate}.  No reduction assumption is needed.
For every prescribed rational cycle $Z$ representing $z$, strong G
holds on exactly $X^{\rm an}\setminus|Z|^{\rm an}$, over the
original $K$, and on every compact actual test.
\end{theorem}

\begin{proof}
Over $\overline K$, algebraically trivial Chow classes are generated
by degree-zero divisor correspondences on smooth projective curves.
This is also a consequence of the curve case of
\cite[Theorems 1--2, Proposition 3.14]{ACMV}.  Write $z_{\overline K}$
as a finite rational sum of such expressions.  All curves,
correspondences, marked points, divisor classes and rational
equivalence witnesses involve finitely many algebraic coefficients.
They therefore descend, after enlarging once, to a finite Galois
extension $L/K$.  Enlarge $L$ also so that each curve has a chosen
$L$-point.  Rational coefficients can either be retained throughout
or cleared by one positive integer at this step.

Lemma~\ref{r12p:poincare}, followed by taking a product of the
finitely many Jacobians and adding the pullbacks of their families,
gives an abelian variety $J/L$, a class $\Theta$ additive on $J(L)$,
and $t\in J(L)$ such that
\[
                              \Theta_t=z_L.
\]
Form the native family $\Theta^N$ from
Lemma~\ref{r12p:norm-family}.  Its fiber at the corresponding
$a\in A(K)$ is
\[
        (\Theta^N)_a=h_*z_L=[L:K]z
                 \quad\text{in }\operatorname{CH}^r(X)_{\mathbb Q},
\tag{R12.P13}\label{r12p:degree-norm}
\]
by the algebraic flat-pullback/proper-pushforward degree formula.
Apply Theorem~\ref{r12p:compact-additive} entirely over $K$ and
divide by $d=[L:K]$ to obtain \eqref{r12p:alg-zero}.

Here is the primitive with every correction exposed.  Choose a
model character $\Omega$ for $\Theta^N$ on $X\times A$, and let
$\beta_a$ be \eqref{r12p:beta}.  If $\Omega_z$ is the prescribed
model character for $z$, then
$\Omega_a-\Omega_0-d\Omega_z$ has generic Chow character zero.
Give it a global primitive $H_z$ by Proposition~\ref{r11p:generic-chow-primitive}.  Then
\[
             \Psi_z=\frac{\beta_a-H_z}{d},
             \qquad D\Psi_z=\Omega_z.
\tag{R12.P14}\label{r12p:final-primitive}
\]
Every term retains \eqref{r12p:certificate}.

Finally take the exact-support pair $(G_Z,\Omega_Z)$ supplied by
the supported-complex construction, and apply the preceding result
to its model curvature.  Put
\[
       \gamma_Z=G_Z-\Psi_Z\big|_{X^{\rm an}\setminus|Z|^{\rm an}}.
\]
The correction is global, so introduces no additional pole.  Stokes
and actual-current injectivity give, for every
$\eta\in\mathcal B_c^{n-r,n-r}(X^{\rm an})$,
\[
 \int_{X^{\rm an}\setminus|Z|^{\rm an}}\gamma_ZD\eta
                         =-\int_{Z^{\rm an}}\eta.
\]
This is the required exact-support, all-actual-test statement.

\end{proof}



% Source: research-sources/green-draft/abelian_numerical.tex
\section{Finite numerical certificates and abelian varieties}
\label{sec:abelian-numerical}

In this section $K/\mathbf Q_p$ is finite, with valuation ring $R$
and finite residue field $k$.  For a smooth integral projective
$K$-variety $V$ of dimension $d$, write
\[
 \CH^r(V)_{\mathbf Q,\homol}
 =\ker\bigl(\CH^r(V)_{\mathbf Q}\longrightarrow
          H^{2r}(V_{\overline K},\mathbf Q_\ell(r))\bigr),
 \qquad \ell\ne p.
\]
When $V$ is not geometrically integral, the geometric class is
understood on every component.  The construction of the preceding
sections applies to smooth integral $V$ with this convention: the
finite multisection argument, model Chern forms, finite trace, and
supported-complex Green construction use integrality over $K$;
geometric integrality enters the algebraic-equivalence theorem only.

\subsection{A finite numerical certificate}

\begin{definition}\label{def:numerical-certificate}
Let $\mathcal V/R$ be integral, projective and flat, with smooth
integral generic fiber $V$.  A virtual model bundle
$\xi\in K_0(\mathcal V)_{\mathbf Q}$ is an
\emph{$r$-numerical certificate} for
$b\in\CH^r(V)_{\mathbf Q}$ if
\[
  \operatorname{ch}_r(\xi|_V)=b
\quad\text{and}\quad
  \deg_k\bigl(\operatorname{ch}_r(\xi_s)\cap[T]\bigr)=0
\]
for every integral $r$-dimensional closed subscheme
$T\subset\mathcal V_s$.  The degree is that of the resulting
zero-cycle on the proper $k$-scheme $T$.
\end{definition}

\begin{lemma}[Model refinements preserve numerical certificates]
\label{lem:refine-certificate}
If $\rho:\mathcal V'\to\mathcal V$ is a projective morphism of
integral projective flat $R$-models that is an isomorphism on
generic fibers, then $\rho^*\xi$ is an $r$-numerical certificate
for the same generic class whenever $\xi$ is.
\end{lemma}

\begin{proof}
The equality on the generic fiber is immediate.  For an integral
vertical $r$-cycle $T'$ on $\mathcal V'$, the operational projection
formula gives
\[
 \deg_k\bigl(\operatorname{ch}_r(\rho^*\xi_s)\cap[T']\bigr)
 =\deg_k\bigl(\operatorname{ch}_r(\xi_s)\cap\rho_{s*}[T']\bigr).
\]
The proper pushforward on the right is zero if the image has
dimension smaller than $r$; otherwise it is a positive integer
multiple of an integral $r$-cycle on $\mathcal V_s$.  The numerical
certificate makes its degree zero.
\end{proof}

\begin{theorem}[Finite numerical descent]\label{thm:finite-numerical}
Let $h:V\to X$ be finite dominant between smooth integral
projective $K$-varieties, and let $1\le r\le\dim X$.  If $h^*z$
has an $r$-numerical certificate on a projective flat model of $V$,
then
\[
                     \Ob^r_{X,K}(z)=0
                     \qquad(z\in\CH^r(X)_{\mathbf Q}).
\]
Each prescribed rational cycle $Z$ in class $z$ consequently has
an ordinary actual zero-curvature Green form on
$X^{\an}\setminus|Z|^{\an}$ satisfying \eqref{eq:main-green}
for every compact actual complementary test.

The assertion also applies when $V$ is the underlying $K$-variety
of a smooth integral projective variety over a finite extension
$F/K$, and its model is projective flat over $F^\circ$.
\end{theorem}

\begin{proof}
Put $e=\deg_K(h)$.  A finite map between smooth integral varieties
of the same dimension is flat, so $h$ is finite flat of rank $e$.
Choose a projective flat model $\mathcal X$ and a virtual bundle
representing $z$; let $\Omega_X$ be its actual degree-$r$ character.
The existence of this model character and its dependence only on
the generic Chow class in actual Bott--Chern cohomology were proved
in Proposition~\ref{r11p:generic-chow-primitive}.

Let $\xi$ be the supplied certificate on $\mathcal V$.  Form the
closure of the generic graph of $h$ in
$\mathcal V\times_R\mathcal X$.  It is a projective flat model
$\mathcal V'$ of $V$, with morphisms to both original models.
By Lemma~\ref{lem:refine-certificate}, the pullback of $\xi$ remains
numerically trivial on vertical $r$-cycles.  Replacing
$\mathcal V$ by this model makes the two model characters
comparable.

The numerical vanishing theorem \Cref{r10c:numerical-zero,r10c:chern}
gives the \emph{strict equality of actual forms}
\[
                     \operatorname{ch}_r(\bar\xi)=0
                                \quad\text{on }V^{\an}.
\]
The virtual bundles $h^*\xi_X$ and $\xi$ have the same generic
degree-$r$ Chow character.  Thus
$h^*\Omega_X-\operatorname{ch}_r(\bar\xi)=h^*\Omega_X$
has a global actual $D$-primitive by the generic-Chow-zero theorem
Proposition~\ref{r11p:generic-chow-primitive}.  We use its constructive form:
\begin{equation}\label{eq:certificate-primitive}
 h^*\Omega_X=\sum_jc_j\,\Omega_j,
 \qquad
 \beta_V=\sum_j\frac{c_j}{e_j}
        \Tr^{\mathscr M}_{b_j}(G_j),
 \qquad
 DG_j=b_j^*\Omega_j.
\end{equation}
Here $c_j\in\mathbf Q$, each $b_j:T_j\to V$ is finite flat of
native degree $e_j$, and $G_j$ is a global sum of the original
PL scalar--curvature generators.  In particular the composite
$h b_j$ is finite flat on the \emph{whole} generic fiber $X$.

The finite trace of \Cref{r11p:minimal-trace-theorem} therefore
defines the following actual form on the original $X^{\an}$:
\begin{equation}\label{eq:descent-primitive}
 \beta_X=\frac1e\sum_j\frac{c_j}{e_j}
                  \Tr^{\mathscr M}_{h b_j}(G_j).
\end{equation}
For an arbitrary compact complementary actual test $\eta$, the
finite-trace pairing and actual Stokes give
\[
 \begin{aligned}
 \int_X\beta_X\wedge D\eta
   &=\frac1e\sum_j\frac{c_j}{e_j}
                   \int_{T_j}G_j\wedge(hb_j)^*D\eta\\
   &=\frac1e\sum_jc_j\int_V\Omega_j\wedge h^*\eta\\
   &=\frac1e\int_Vh^*(\Omega_X\wedge\eta)
     =\int_X\Omega_X\wedge\eta.
 \end{aligned}
\]
The last equality uses the full native degree $e$, and the first
three equalities hold for \emph{every} compact actual test by the
finite-trace theorem.  Mihatsch's injectivity of the map from
actual forms to currents now yields $D\beta_X=\Omega_X$.
The obstruction class vanishes; subtracting $\beta_X$ from the
exact-support curved pair of Corollary~\ref{r11p:any-cycle-pair} yields
the asserted Green form.

For a model over $F^\circ$, regard it as a model over $R$.
The scheme-theoretic special fiber retains its ramification
multiplicity.  Each vertical numerical degree over $k$ is the
corresponding degree over the residue field of $F$ multiplied
by that residue-field degree.  The same proof uses the complete
native generic degree of $h$, including the constant-field factor.
\end{proof}

\subsection{The special fiber of a semistable abelian model}

Let $R'$ be a complete discrete valuation ring with finite residue
field $k'$ and fraction field $F$.  Let $\mathcal A/R'$ be a regular
projective strictly semistable model of an abelian variety $A/F$
of dimension $d$, and let $Y=\bigcup_iD_i$ be its reduced special
fiber.  With $Y^{(1)}=\coprod_iD_i$ and
$Y^{(2)}=\coprod_{i<j}(D_i\cap D_j)$, define
\[
 A^s(D)=\operatorname{im}\bigl(\CH^s(D)_{\mathbf Q}
           \to H^{2s}(D_{\bar k'},\mathbf Q_\ell(s))\bigr)
\]
for each smooth projective stratum $D$ and set
\begin{equation}\label{eq:kun-strata}
 A^s(Y)=\ker\bigl(A^s(Y^{(1)})\xrightarrow{\rho}A^s(Y^{(2)})\bigr),
 \qquad
 A_s(Y)=\operatorname{coker}\bigl(
 A^{d-s-1}(Y^{(2)})\xrightarrow{\gamma}A^{d-s}(Y^{(1)})\bigr).
\end{equation}
The two arrows are the signed restrictions and Gysin maps.
The natural map $\CH_s(Y)_{\mathbf Q}\twoheadrightarrow A_s(Y)$
is surjective.

\begin{theorem}[K\"unnemann's local semistable theorem]
\label{thm:kunnemann-local}
After a finite extension of any finite $p$-adic field $K$, every
abelian variety over $K$ admits a regular projective strictly
semistable model $\mathcal A/R'$ for which, in all degrees, the
cycle-class map
\[
 A_s(Y)\longrightarrow
 H^{\mathrm{\acute et}}_{2s}(Y_{\bar k'},\mathbf Q_\ell(-s))
\]
is injective, and the middle term of
\begin{equation}\label{eq:kun-complex}
 A_{d+1-r}(Y)\xrightarrow{i^*i_*} A^r(Y)
                   \xrightarrow{\cap[Y]} A_{d-r}(Y)
\end{equation}
is exact, where $i:Y\hookrightarrow\mathcal A$.
\end{theorem}

\begin{proof}[Source of the published theorem]
This is the complete-DVR construction and special-fiber theorem
of \cite[Theorems~5.6 and~6.6]{Kunnemann}.
The base ring is complete and has finite residue field, as in
Section~6.2 of the cited work.  The groups in
\eqref{eq:kun-strata} are defined using the geometric
cycle-class images on the individual strata.
\end{proof}

\begin{lemma}[A pure, vertically numerical extension]
\label{lem:abelian-lift}
Let $A/F$ have a model furnished by
\Cref{thm:kunnemann-local}.  For every geometrically homologically
trivial class $z\in\CH^r(A)_{\mathbf Q}$,
$1\le r\le d$, there is $\Theta_0\in\CH^r(\mathcal A)_{\mathbf Q}$
restricting to $z$, such that its refined restriction to
every $D_i$ is homologically trivial.  In particular $z$ has
an $r$-numerical certificate on $\mathcal A$.
\end{lemma}

\begin{proof}
Choose a Chow extension $\Theta$ of $z$ by taking closures of
cycle representatives.  Its tuple $a$ of component restriction
classes belongs to $A^r(Y)$.  The class of
$a\cap[Y]\in A_{d-r}(Y)$ is the specialization of $z$.
Compatibility of proper specialization with geometric
$\ell$-adic cycle classes makes its image in geometric homology
zero.  By the injectivity in \Cref{thm:kunnemann-local},
$a\cap[Y]=0$.

Exactness of \eqref{eq:kun-complex} gives
$b\in A_{d+1-r}(Y)$ with $i^*i_*b=a$.
Choose $\beta\in\CH_{d+1-r}(Y)_{\mathbf Q}$ mapping to $b$
and set
\[
                           \Theta_0=\Theta-i_*\beta.
\]
The generic fiber of $i_*\beta$ is empty, so $\Theta_0|_A=z$.
The tuple of its component restrictions is
$a-i^*i_*b=0$ in $A^r(Y)$.  Since $A^r(Y)$ is a subspace of
$\bigoplus_i A^r(D_i)$, each component restriction has zero
geometric $\ell$-adic class and is numerically trivial.

The regularity of $\mathcal A$ and the rational Chern-character
isomorphism on regular schemes
\cite[Tag 0FEY]{StacksFEY} supply a unique
$\xi\in K_0(\mathcal A)_{\mathbf Q}$ satisfying
\begin{equation}\label{eq:pure-ch}
                    \operatorname{ch}(\xi)\cap[\mathcal A]
                              =\Theta_0.
\end{equation}
For completeness, applying the Adams operation $2^{-r}\psi^2$
to \eqref{eq:pure-ch} yields the same codimension-$r$ class.
Uniqueness therefore gives $\psi^2\xi=2^r\xi$.
Because $\operatorname{ch}_j(\psi^2\xi)=2^j\operatorname{ch}_j(\xi)$,
all operational character components except $j=r$ vanish.

Every integral vertical $r$-cycle $T\subset\mathcal A_s$
is contained in a component $D_i$.  The projection formula for
the refined restriction gives
\[
 \deg_{k'}\bigl(\operatorname{ch}_r(\xi_s)\cap[T]\bigr)
 =\deg_{k'}\bigl((\Theta_0|_{D_i})\cap[T]\bigr)=0.
\]
The last equality follows from numerical triviality of the
component restriction.  Thus $\xi$ is a numerical certificate.
\end{proof}

\subsection{Zero-curvature Green forms on abelian targets}

\begin{theorem}[Full actual strong Green forms for abelian varieties]
\label{thm:abelian-G}
Let $A/K$ be an abelian variety of arbitrary dimension.  For
every $1\le r\le\dim A$ and every prescribed rational cycle $Z$
whose Chow class is geometrically homologically trivial, there is
an actual form
\[
 \gamma_Z\in\B^{r-1,r-1}
                    (A^{\an}\setminus|Z|^{\an})
\]
satisfying \eqref{eq:main-green} for all compact actual tests
over the original $K$.
\end{theorem}

\begin{proof}
Choose a finite field extension $F/K$ for which $A_F$ has the
semistable model of \Cref{thm:kunnemann-local}.  Its $r$-numerical
certificate for $Z_F$ is supplied by
Lemma~\ref{lem:abelian-lift}.  The finite dominant map
$h:A_F\to A$ has native degree $[F:K]$ and is finite flat as
a map of smooth integral $K$-varieties.  Apply
\Cref{thm:finite-numerical} to $h$.  It supplies a global actual
primitive for the obstruction of $Z$ on the original
$A^{\an}$, and hence the asserted Green form on the prescribed
support complement.
\end{proof}

\begin{corollary}[Finite abelian domination]
\label{cor:finite-abelian}
Let $X/K$ be smooth projective geometrically integral, and assume
there is a finite surjective morphism $f:A\to X$ from an abelian
variety $A/K$.  Every geometrically homologically trivial rational
cycle on $X$ satisfies the strong Green identity in all
codimensions, with prescribed exact support and all compact
original-field actual tests.
\end{corollary}

\begin{proof}
For $z\in\CH^r(X)_{\mathbf Q,\homol}$, the class
$f^*z$ is geometrically homologically trivial by functoriality
of the geometric cycle-class map.  Extend $K$ to $F$ as in
\Cref{thm:abelian-G}; Lemma~\ref{lem:abelian-lift} gives a
numerical certificate for its pullback to $A_F$.
The composite finite dominant map $A_F\to A\to X$ has
degree $[F:K]\deg f$.  Apply \Cref{thm:finite-numerical}
directly to this composite.
\end{proof}



% Source: research-sources/green-draft/chow_correspondence.tex
\section{Numerical certificates under Chow correspondences}
\label{sec:correspondences}

This section uses ordinary rational intersection theory on regular
projective models.  The only analytic input is the finite numerical
descent theorem already proved in \Cref{thm:finite-numerical}.

\subsection{A supported numerical projection formula}

Let $\mathcal V/R$ be a regular, integral, projective and flat model,
and let $\Theta\in\CH^p(\mathcal V)_{\mathbf Q}$ be homogeneous.
By the rational Chern-character isomorphism for regular schemes
\cite[Tag~0FEY]{StacksFEY}, there is a unique
$\xi_\Theta\in K_0(\mathcal V)_{\mathbf Q}$ such that
\begin{equation}\label{eq:op-pure-character}
       \operatorname{ch}(\xi_\Theta)\cap[\mathcal V]=\Theta.
\end{equation}
The Adams identity, applied to \eqref{eq:op-pure-character}, gives
$\psi^2\xi_\Theta=2^p\xi_\Theta$ and hence
$\operatorname{ch}_j(\xi_\Theta)=0$ as an operational class
for $j\ne p$.  We write
$c_\Theta=\operatorname{ch}_p(\xi_\Theta)$.  It is an
operational class on every closed support and satisfies
$c_\Theta\cap[\mathcal V]=\Theta$.

We call $\Theta$ \emph{vertically $p$-numerically trivial} if
\begin{equation}\label{eq:vertical-numerical-def}
 \deg_k(c_\Theta\cap[T])=0
 \quad\text{for every integral vertical $p$-cycle}
 \quad T\subset\mathcal V_s.
\end{equation}
Equivalently, $\xi_\Theta$ is a $p$-numerical certificate for
$\Theta$ on the generic fiber.

\begin{lemma}[Numerical pullback]\label{lem:operational-pullback}
Let $f:\mathcal W\to\mathcal V$ be a projective morphism between
regular integral projective flat $R$-models.  If $\Theta$ is
vertically $p$-numerically trivial, then its refined pullback
$f^!\Theta=(f^*c_\Theta)\cap[\mathcal W]$ is vertically
$p$-numerically trivial.
\end{lemma}

\begin{proof}
The schemes are regular, so $f$ is a local complete intersection
morphism, and the projective factorization supplies its refined
Gysin operation.  For any integral vertical $p$-cycle
$T\subset\mathcal W_s$, the proper projection formula gives
\[
 \deg_k((f^*c_\Theta)\cap[T])
       =\deg_k(c_\Theta\cap f_*[T])=0.
\]
Indeed $f_*[T]$ is a rational linear combination of vertical
$p$-cycles, or zero.  The displayed refined pullback has pure
operational representative $f^*c_\Theta$ by
\eqref{eq:op-pure-character} and Gysin compatibility.
\end{proof}

\begin{lemma}[Supported numerical projection]
\label{lem:supported-numerical-projection}
Let $f:\mathcal W\to\mathcal V$ be a projective morphism between
regular integral projective flat $R$-models of total dimensions
$M$ and $N$.  Let $\Theta\in\CH^p(\mathcal W)_{\mathbf Q}$ be
vertically $p$-numerically trivial.  Put
$\delta=M-N$, $q=p-\delta$, and assume $0\le q\le N-1$.
Then $f_*\Theta\in\CH^q(\mathcal V)_{\mathbf Q}$ is vertically
$q$-numerically trivial.
More precisely, for every integral vertical $q$-cycle
$T\subset\mathcal V_s$ one has the supported projection identity
\begin{equation}\label{eq:supported-identity}
 (f_*\Theta)\cdot[T]
   = f_*\bigl(\Theta\cdot f^![T]\bigr)
                      \qquad\text{in }\CH_0(\mathcal V)_{\mathbf Q}.
\end{equation}
\end{lemma}

\begin{proof}
Since $\mathcal V$ and $\mathcal W$ are regular, the projective
morphism $f$ is a local complete intersection morphism: factor it
through a projective bundle over $\mathcal V$, then use the
regular embedding of the regular source into that regular bundle.
The fundamental class $[T]$ has codimension $N-q$ in the regular
scheme $\mathcal V$.  Its refined pullback
\[
       f^![T]\in\CH_{M-(N-q)}(\mathcal W)_{\mathbf Q}
                   =\CH_p(\mathcal W)_{\mathbf Q}
\]
is supported on $f^{-1}(T)$, hence is a rational combination of
vertical $p$-cycles.  The bivariant projection formula for the
lci pullback on regular schemes
\cite[Chapter~17]{Fulton} and \cite[Tags~0FF6 and 0FFA]{StacksGysin}
gives \eqref{eq:supported-identity}.
The degree of $\Theta\cdot f^![T]$ is zero by the numerical
hypothesis on $\Theta$.  Since $c_{f_*\Theta}\cap[T]$ is the
intersection product $(f_*\Theta)\cdot[T]$, the degree required
for vertical numerical triviality vanishes.
\end{proof}

\begin{lemma}[Products of numerical and operational classes]
\label{lem:numerical-times-class}
Let $\mathcal W$ be a regular projective flat model.  Suppose
$\alpha\in\CH^p(\mathcal W)_{\mathbf Q}$ is vertically
$p$-numerically trivial and
$\beta\in\CH^t(\mathcal W)_{\mathbf Q}$.
If $p+t$ lies between zero and the relative dimension of
$\mathcal W$, then $\alpha\cdot\beta$ is vertically
$(p+t)$-numerically trivial.
\end{lemma}

\begin{proof}
For an integral vertical $(p+t)$-cycle $T$, the operational
intersection $c_\beta\cap[T]$ is a rational combination of
vertical $p$-cycles.  The degrees of
$c_\alpha\cap(c_\beta\cap[T])$ vanish termwise.  Purity of
the rational operational representatives identifies
$c_{\alpha\beta}=c_\alpha c_\beta$.
\end{proof}

\subsection{Transfer along algebraic correspondences}

Let $V$ and $X$ be smooth projective geometrically integral
$F$-varieties of dimensions $m$ and $n$.  A correspondence
$\Gamma\in\CH^{m+a}(V\times_F X)_{\mathbf Q}$ acts by
\[
 \Gamma_*b=(\operatorname{pr}_X)_*
             (\operatorname{pr}_V^*b\cdot\Gamma)
       \in\CH^{r+a}(X)_{\mathbf Q},
       \qquad b\in\CH^r(V)_{\mathbf Q}.
\]

\begin{proposition}[Algebraic correspondence transfer]
\label{prop:correspondence-transfer}
Assume that $V$ and $X$ have regular integral projective flat
$R_F$-models $\mathcal V$ and $\mathcal X$.  Suppose that
$b\in\CH^r(V)_{\mathbf Q}$ has an $r$-numerical certificate
on $\mathcal V$.  For every correspondence
$\Gamma\in\CH^{m+a}(V\times X)_{\mathbf Q}$ with
$1\le r+a\le n$, its image $\Gamma_*b$ has an
$(r+a)$-numerical certificate on $\mathcal X$.
\end{proposition}

\begin{proof}
Write $R'=R_F$.  The projective flat product
$\mathcal V\times_{R'}\mathcal X$ is integral: its generic
fiber $V\times_F X$ is geometrically integral, and flatness
excludes any additional irreducible component supported on the
special fiber.  By the projective semistable alteration theorem
of de Jong \cite[Theorem~6.5]{deJong}, there is a projective
generically finite morphism
\[
 g:\mathcal W\longrightarrow
                  \mathcal V\times_{R'}\mathcal X
\]
from a regular integral projective flat model (possibly over a
finite extension of $R'$).  The native degree of its generic
map $g_F:W\to V\times_F X$ is a positive integer $d$.
Indeed, de Jong first constructs a semistable open source inside
a projective compactification.  Since the target here is proper,
the alteration source is proper over the base, so this open
immersion is closed as well.  The integral compactification
therefore equals the source.
Let $f_V$ and $f_X$ be the composite model morphisms to the
respective factors.

Choose a pure vertically numerical class
$\Theta\in\CH^r(\mathcal V)_{\mathbf Q}$ with
$\Theta|_V=b$.  The generic class $g_F^*\Gamma$ has a Chow
extension $\widetilde\Gamma\in\CH^{m+a}(\mathcal W)_{\mathbf Q}$;
the closure of cycle representatives gives such an extension.
Set
\[
 \Xi=(f_V^!\Theta)\cdot\widetilde\Gamma
      \in\CH^{r+m+a}(\mathcal W)_{\mathbf Q}.
\]
By Lemmas~\ref{lem:operational-pullback}
and~\ref{lem:numerical-times-class},
$\Xi$ is vertically $(r+m+a)$-numerically trivial.  Since
$\dim\mathcal W-\dim\mathcal X=m$, the supported numerical
projection lemma applies to $f_X$, and gives the vertically
$(r+a)$-numerically trivial class
\[
          \Lambda=\frac1d(f_X)_*\Xi
                     \in\CH^{r+a}(\mathcal X)_{\mathbf Q}.
\]

On smooth generic fibers, pullback is an ordinary lci Chow
pullback, and the projection formula for the generically finite
$g_F$ gives
\[
 \begin{aligned}
 \Lambda|_X
  &=\frac1d(\operatorname{pr}_X)_*
       (g_F)_*g_F^*(\operatorname{pr}_V^*b\cdot\Gamma)\\
  &=(\operatorname{pr}_X)_*
                      (\operatorname{pr}_V^*b\cdot\Gamma)
   =\Gamma_*b.
 \end{aligned}
\]
The pure virtual model bundle $\xi_\Lambda$ defined by
\eqref{eq:op-pure-character} is the required numerical
certificate.
\end{proof}

\begin{theorem}[A Chow-correspondence criterion for actual Green forms]
\label{thm:chow-abelian-criterion}
Let $X/K$ be smooth projective geometrically integral.  Suppose
there exists a finite extension $F/K$ for which the following
properties hold:
\begin{enumerate}
\item $X_F$ has a regular integral projective flat $R_F$-model;
\item there is an abelian variety $B/F$ with a regular projective
  semistable model satisfying \Cref{thm:kunnemann-local};
\item there are finitely many correspondences
  $u_j:X_F\rightsquigarrow B$ and
  $v_j:B\rightsquigarrow X_F$ of opposite degree shifts
  with
  \begin{equation}\label{eq:abelian-diagonal}
                \sum_jv_j\circ u_j=\Delta_{X_F}
       \qquad\text{in }\CH^{\dim X}(X_F\times X_F)_{\mathbf Q}.
  \end{equation}
\end{enumerate}
Then every prescribed rational cycle $Z$ on $X$ with
geometrically homologically trivial class admits an actual
zero-curvature Green form on the exact complement
$X^{\an}\setminus |Z|^{\an}$ satisfying
\eqref{eq:main-green} against every compact actual test over $K$.
\end{theorem}

\begin{proof}
Write $z=[Z]$ and $z_F=[Z_F]$.  Equation
\eqref{eq:abelian-diagonal} gives
\[
             z_F=\sum_j(v_j)_*(u_j)_*z_F.
\]
Each $(u_j)_*z_F$ is geometrically homologically trivial, by
compatibility of algebraic correspondences with geometric
$\ell$-adic cycle classes.  Terms of codimension outside
$1,\ldots,\dim B$ vanish (in codimension zero this follows
from homological triviality); for the remaining terms
Lemma~\ref{lem:abelian-lift} supplies a numerical certificate on
the regular semistable model of $B$.

Apply Proposition~\ref{prop:correspondence-transfer} to $v_j$, taking the
given regular model of $X_F$ as the target.  Its resulting
numerical certificates add to a certificate for $z_F$.
Finally, $X_F\to X$ is finite flat of native degree $[F:K]$.
The finite numerical descent theorem \Cref{thm:finite-numerical}
produces the asserted actual Green form on the original $X$.
\end{proof}

\subsection{Products of curves and abelian varieties}

\begin{lemma}[Chow diagonal of a product of curves]
\label{lem:curve-product-diagonal}
Let $F$ have characteristic zero, let $A/F$ be an abelian variety,
and let $C_1,\ldots,C_s/F$ be smooth projective geometrically
integral curves with $F$-rational points.  Put
$X=A\times_F C_1\times_F\cdots\times_F C_s$.
Then the diagonal of $X$ admits a finite decomposition
\eqref{eq:abelian-diagonal}, where every intermediate variety
is the single abelian variety
$B=A\times_F\prod_i\operatorname{Jac}(C_i)$,
and all correspondences are defined over $F$.
\end{lemma}

\begin{proof}
Choose points $P_i\in C_i(F)$.  The three mutually orthogonal
Chow projectors
\[
 \pi_{0,i}=[P_i\times C_i],\qquad
 \pi_{2,i}=[C_i\times P_i],\qquad
 \pi_{1,i}=\Delta_{C_i}-\pi_{0,i}-\pi_{2,i}
\]
split the motive of each curve into $\mathbf 1$, its
degree-one motive, and the Lefschetz motive $\mathbb L$.
Idempotence follows from $\deg(P_i)=1$ and the mixed
compositions vanish because $[P_i]^2=0$ on the curve.
The Abel--Jacobi map and the Poincar\'e correspondence identify
the degree-one motive with the degree-one direct summand of
the Chow motive of $J_i=\operatorname{Jac}(C_i)$.  The latter
is split by the rational Chow--K\"unneth projector of the
abelian variety, so the inverse correspondences exist in the
category of rational Chow motives.  These standard identities
are established in \cite[Sections~3 and~5]{Scholl}.

Take the tensor product of the finitely many splittings, including
the identity motive of $A$.  Every resulting summand factors
through the motive of $A$ times a product of some $J_i$, with
an integral Tate twist recording the number of Lefschetz factors.
The projection $\pi$ from $B=A\times\prod_i J_i$ onto any such
subproduct admits a zero-section $i$.  The identity
$i^*\pi^*=\operatorname{id}$, implemented by transpose-graph
correspondences, splits the motive of the subproduct from
the motive of $B$.  Thus every summand can be
factored through this one abelian variety, with the same Tate
shift.
Express the inclusion and retraction morphisms as algebraic
correspondences $u_j$ and $v_j$; their composition is the tensor
product of the chosen projectors.  Summing over all choices of
the three projectors on each curve gives
$\sum_jv_j\circ u_j=\Delta_X$ in rational Chow groups.
\end{proof}

\begin{lemma}[Regular models for products]
\label{lem:product-model}
Let $A/K$ be an abelian variety and let $C_1,\ldots,C_s/K$
be smooth projective geometrically integral curves.  After
one finite extension $F/K$ there are simultaneously:
\begin{enumerate}
\item rational points on all the $C_{i,F}$;
\item regular projective strictly semistable models of the
  $C_{i,F}$, and a regular projective strictly semistable
  model of $A_F\times\prod_i C_{i,F}$;
\item a regular projective semistable model of the abelian
  variety $B=A_F\times\prod_i\operatorname{Jac}(C_{i,F})$
  satisfying \Cref{thm:kunnemann-local}.
\end{enumerate}
\end{lemma}

\begin{proof}
Each curve acquires a rational point and a projective regular
strictly semistable model after finite extension, by the
semistable reduction theorem for curves.  Enlarge this initial
field so that $A$ has a regular projective strictly semistable
model by \Cref{thm:kunnemann-local}.  Now apply that theorem to
the single abelian variety
$B=A\times\prod_i\operatorname{Jac}(C_i)$, obtaining a
possibly larger field on which $B$ has the required model.
The earlier models of $A$ and the curves are carried to this
final field by the ramified base-change regularization below.
Semistable
curve and abelian models remain regular and strictly semistable
after an unramified extension.  After ramified extension,
Hartl's explicit sequence of blowups recovers such a model
without changing the generic fiber
\cite[Proposition~2.2]{Hartl}.  Finally, his regularization of
the product of two strict semistable models
\cite[Proposition~2.1]{Hartl} can be iterated over the finite
family.  All operations are projective blowups of projective
schemes; thus the resulting model is projective, regular and
flat, with the prescribed smooth generic fiber.
\end{proof}

\begin{theorem}[Products of curves and abelian varieties]
\label{thm:curve-products-G}
Let $A/K$ be an abelian variety and let $C_1,\ldots,C_s/K$
be smooth projective geometrically integral curves, with
$K/\mathbf Q_p$ finite.  On
\[
                       X=A\times_K C_1\times_K\cdots\times_K C_s,
\]
every prescribed rational cycle of codimension
$1\le r\le\dim X$ whose class is geometrically homologically
trivial has an actual zero-curvature Green form satisfying
\eqref{eq:main-green} on its exact support complement over the
original field, against all compact actual tests.
\end{theorem}

\begin{proof}
Choose $F/K$ and the regular semistable models of
Lemma~\ref{lem:product-model}.  The Chow-diagonal identity of
Lemma~\ref{lem:curve-product-diagonal} meets every hypothesis of
\Cref{thm:chow-abelian-criterion}.  Apply that theorem.
\end{proof}

\begin{corollary}[Finite domination by products]
\label{cor:finite-curve-product}
Let $Y/K$ be smooth projective geometrically integral, and
suppose that there is a finite surjective morphism
$A\times\prod_i C_i\to Y$ from a product as in
\Cref{thm:curve-products-G}.  Then the conclusion of that
theorem holds for every geometrically homologically trivial
rational cycle on $Y$.
\end{corollary}

\begin{proof}
Pull a cycle back through the finite map and choose the common
extension and models of Lemma~\ref{lem:product-model}.  The proof
of \Cref{thm:chow-abelian-criterion} constructs an actual
numerical certificate for the pullback on a regular model of
the product after that extension.  The composite map from
the base-changed product to $Y$ is finite flat and has degree
$[F:K]$ times the degree of the original finite map.  Apply
\Cref{thm:finite-numerical} directly to this composite.
\end{proof}



\section{The conditional geometric conclusions}\label{sec:conditional-conclusions}

\subsection{Algebraic equivalence over the original field}\label{sec:conditional-algebraic-proof}

\begin{proof}[Proof of \Cref{thm:main-algebraic}]
Write $z=[Z]$. Choose a finite Galois extension $L/K$ over which the curves, degree-zero divisors and correspondence identities expressing $z_L$ by algebraic equivalence are defined. \Cref{r12p:poincare} provides an additive family $\Theta$ on $X_L\times J$, for a product of Jacobians $J/L$, with a fibre $\Theta_t=z_L$.

Let $A=\operatorname{Res}_{L/K}J$. Evaluation $A_L\to J$ and algebraic finite-flat pushforward give the norm family of \Cref{r12p:norm-family}. It is additive and has a $K$-rational fibre
\[
 (\Theta^N)_a=[L:K]z.
\]
By (A1), the family defines the additive obstruction map
\[
 \varphi:A(K)\longrightarrow H^{r,r}_{\BC}(X),
 \qquad b\longmapsto\Ob^r_{X,K}((\Theta^N)_b).
\]
Represent the family obstruction by a global model Chern form. Property (R) makes its pulled-back form constant near zero, so $\ker\varphi$ contains an identity neighbourhood. Compactness of $A(K)$ implies that a positive multiple of any point lies in that neighbourhood, as in \Cref{r12p:compact-multiple}. The target is torsion-free, hence $\varphi=0$. Thus $[L:K]\Ob^r_{X,K}(z)=0$, and division gives $\Ob^r_{X,K}(z)=0$.

Use (A2) for the original representative $Z$, not a moved cycle. Since $[\Omega_Z]=0$, choose a global actual $\Psi_Z$ with $D\Psi_Z=\Omega_Z$. Set $\gamma_Z=G_Z-\Psi_Z|_{U_Z}$. Actual Stokes and \eqref{eq:main-pair} give \eqref{eq:main-green} for every stated test, on the same open set $U_Z$.
\end{proof}

\subsection{Abelian correction and diagonal transfer}\label{sec:conditional-abelian-proof}

\begin{proof}[Proof of \Cref{thm:main-abelian}]
First let $A$ be an abelian variety over a finite extension of $K$. Choose a further finite extension with the semistable model of \Cref{thm:kunnemann-local}. For a geometrically homologically trivial class, \Cref{lem:abelian-lift} gives a pure virtual model bundle whose special-fibre character is numerically trivial on all vertical cycles of the required dimension. Property (N), applied to the native finite map to the original $K$-variety, gives zero obstruction. This uses the entire degree of the extension.

Now fix the retraction data \eqref{eq:construction-retraction} and $z\in\CH^r(X)_{\mathbf Q,\homol}$. Every $(u_j)_*z_L$ is geometrically homologically trivial by compatibility of correspondences with cycle classes. Terms outside positive codimension on the abelian variety vanish; in codimension zero this follows from homological triviality. The first paragraph therefore annihilates the obstruction of each remaining term. Property (C) annihilates the obstruction of its image under $v_j$. Summing the Chow identity gives zero obstruction for $z_L$.

Apply the algebraic pushforward $\pi_*:X_L\to X$ and (C). Since $\pi_*z_L=[L:K]z$, rational linearity gives $\Ob^r_{X,K}(z)=0$. If the correspondences act as $\lambda\id$ rather than the identity in the chosen codimension, the same argument gives $\lambda[L:K]\Ob^r_{X,K}(z)=0$, and $\lambda\ne0$ permits division. Finally apply (A2) and the global correction used in the preceding proof.
\end{proof}

\begin{thebibliography}{99}
\bibitem{M} A. Mihatsch, \emph{$\delta$-Forms on Lubin--Tate spaces}, arXiv:2112.10018, version 4 (2024), \url{https://arxiv.org/abs/2112.10018v4}.
\bibitem{MTrop} A. Mihatsch, \emph{On tropical intersection theory}, Selecta Math. (N.S.) \textbf{29} (2023), art. 17, \url{https://doi.org/10.1007/s00029-022-00818-7}.
\bibitem{CLD} A. Chambert-Loir and A. Ducros, \emph{Formes diff\'erentielles r\'eelles et courants sur les espaces de Berkovich}, arXiv:1204.6277, version 2 (2025), \url{https://arxiv.org/abs/1204.6277}.
\bibitem{GM} W. Gubler and F. Martin, \emph{On Zhang's semipositive metrics}, Doc. Math. \textbf{24} (2019), 331--372, \url{https://doi.org/10.4171/DM/683}.
\bibitem{Fulton} W. Fulton, \emph{Intersection Theory}, second edition, Ergebnisse der Mathematik, Springer, 1998.
\bibitem{Kunnemann} K. K\"unnemann, \emph{Height pairings for algebraic cycles on abelian varieties}, Ann. Sci. \'Ecole Norm. Sup. (4) \textbf{34} (2001), 503--523, \url{https://www.numdam.org/article/ASENS_2001_4_34_4_503_0.pdf}.
\bibitem{StacksFEY} The Stacks Project Authors, \emph{Lemma 42.58.1}, Tag 0FEY, \url{https://stacks.math.columbia.edu/tag/0FEY}.
\bibitem{StacksFEW} The Stacks Project Authors, \emph{Proposition 42.57.1}, Tag 0FEW, \url{https://stacks.math.columbia.edu/tag/0FEW}.
\bibitem{StacksGysin} The Stacks Project Authors, \emph{Lemmas 42.59.7 and 42.59.13}, Tags 0FF6 and 0FFA, \url{https://stacks.math.columbia.edu/tag/0FF6}, \url{https://stacks.math.columbia.edu/tag/0FFA}.
\bibitem{deJong} A. J. de Jong, \emph{Smoothness, semi-stability and alterations}, Publ. Math. Inst. Hautes \'Etudes Sci. \textbf{83} (1996), 51--93, \url{https://www.numdam.org/item/PMIHES_1996__83__51_0.pdf}.
\bibitem{Hartl} U. T. Hartl, \emph{Semi-stability and base change}, Arch. Math. (Basel) \textbf{77} (2001), 215--221, \url{https://doi.org/10.1007/PL00000484}.
\bibitem{Scholl} A. J. Scholl, \emph{Classical motives}, in \emph{Motives}, Proc. Sympos. Pure Math. \textbf{55}, Part 1 (1994), 163--187, \url{https://www.dpmms.cam.ac.uk/~ajs1005/preprints/classical_motives.pdf}.
\bibitem{Stacks} The Stacks Project Authors, \emph{Stacks Project}, particularly Tags 0FB3, 0FBE, 0FET, 0FEW, 0D4I, 0FD6, 0G4C and 00R4, \url{https://stacks.math.columbia.edu}.
\bibitem{MilneJV} J. S. Milne, \emph{Jacobian varieties}, in \emph{Arithmetic Geometry}, G. Cornell and J. H. Silverman (eds.), Springer, 1986, \url{https://www.jmilne.org/math/xnotes/JVs.pdf}.
\bibitem{ACMV} J. D. Achter, S. Casalaina-Martin and C. Vial, \emph{Parameter spaces for algebraic equivalence}, Int. Math. Res. Not. IMRN (2019), no. 6, 1863--1893.
\end{thebibliography}
\end{document}
